% La Philosophie africaine — Philosophie Tle Tech — Cours signé OnBuch+
% Programme officiel MINESEC. Compiler avec : tectonic N01-philosophie-africaine.tex
\newcommand{\DOCMATIERE}{Philosophie}
\newcommand{\DOCNIVEAU}{Tle Tech}
\newcommand{\DOCMODULE}{Philosophie – classe de Terminale technique}
\newcommand{\DOCLECON}{N01}
\newcommand{\DOCTITRE}{La Philosophie africaine}
% OnBuch+ — préambule « cours signés » ENRICHI (compilé avec tectonic / XeLaTeX)
% Système visuel « Pop » d'OnBuch+ : fond crème (jamais blanc), bordures encre
% épaisses, ombres « dures » décalées, titres Archivo Black en capitales.
% Version MATHÉMATIQUES : graphes (pgfplots/tikz), boîtes pédagogiques
% (définition, propriété, méthode, attention, le savais-tu, exercice résolu,
% exercice, TP, bilan), page de garde et signature OnBuch+ ancrée en bas.
%
% Macros attendues dans main.tex AVANT \input{preamble} :
%   \DOCMATIERE  \DOCNIVEAU  \DOCMODULE  \DOCLECON (numéro)  \DOCTITRE
\documentclass[11pt]{article}

\usepackage{fontspec}
\usepackage[french]{babel}
\usepackage[a4paper,margin=2cm,top=3.4cm,bottom=2.8cm,headheight=1.4cm,footskip=1.3cm]{geometry}
\usepackage{amsmath,amssymb,amsfonts}
\usepackage{mathtools} % \xmapsto, \coloneqq... souvent utilisés par les modèles
\usepackage{cancel} % simplifications barrées (bilans de forces)
% \abs{z} (module, valeur absolue) et \norm{u} (norme), utilisés par les modèles
\providecommand{\abs}{}\let\abs\relax\DeclarePairedDelimiter{\abs}{\lvert}{\rvert}
\providecommand{\norm}{}\let\norm\relax\DeclarePairedDelimiter{\norm}{\lVert}{\rVert}
\usepackage[table]{xcolor} % [table] : fournit aussi \rowcolors (alternance de lignes)
\usepackage{pagecolor}
\usepackage{tikz}
\usetikzlibrary{arrows.meta,positioning,calc,shapes.geometric,shapes.misc,shapes.arrows,decorations.pathmorphing,decorations.pathreplacing,decorations.markings,patterns,shadows,mindmap,trees,fit,backgrounds,matrix,intersections,angles,quotes}
\usepackage{pgfplots}
\usepackage[version=4]{mhchem} % équations nucléaires \ce{^{235}_{92}U}
\usepackage[european,siunitx]{circuitikz} % schémas électriques
\pgfplotsset{compat=1.18}
\usepgfplotslibrary{fillbetween,groupplots} % aires entre courbes (calcul intégral)
\usepackage{enumitem}
\usepackage{fancyhdr}
% [most] charge toutes les bibliothèques tcolorbox (listings, minted, raster,
% documentation...) et gonfle inutilement la mémoire TeX sur un document long
% (~300 boîtes) : on ne charge que ce qui est réellement utilisé.
\usepackage[skins,breakable]{tcolorbox}
\usepackage{graphicx}
\usepackage{colortbl}
\usepackage{booktabs}
\usepackage{tabularx}
\usepackage{multirow}
\usepackage{diagbox} % case d'en-tête diagonale des tableaux à double entrée
% Colonnes extensibles centrée / gauche / droite (C, L, R), souvent utilisées par les modèles
\newcolumntype{C}{>{\centering\arraybackslash}X}
\newcolumntype{L}{>{\raggedright\arraybackslash}X}
\newcolumntype{R}{>{\raggedleft\arraybackslash}X}
\usepackage{array}
\usepackage{multicol}
\usepackage{siunitx}
% parse-numbers=false : accepte aussi \SI{2,5 \times 10^{-3}}{...}
\sisetup{locale=FR,output-decimal-marker={,},group-minimum-digits=4,parse-numbers=false}
\DeclareSIUnit\ohm{\ensuremath{\symup{\Omega}}} % U+2126 absent de la police
\DeclareSIUnit\division{div} % réglages d'oscilloscope (ms/div, V/div)
\DeclareSIUnit\tr{tr} % tours (vitesse de rotation, ex. \SI{1500}{\tr\per\minute})
\DeclareSIUnit\molar{M} % concentration molaire (ex. \SI{3}{\molar}, \SI{1}{\micro\molar})
\DeclareSIUnit\an{an} % année (le modèle confond parfois avec \year, un compteur TeX) — ex. \SI{5}{\kilo\meter\per\an}
\DeclareSIUnit\FCFA{FCFA} % franc CFA (ex. \SI{500}{\FCFA\per\kilo\gram})
\DeclareSIUnit\degreeBrix{\textdegree Brix} % teneur en sucre d'un sirop/jus (réfractométrie)
\DeclareSIUnit\month{mois} % le modèle utilise parfois \month, un compteur TeX, comme unité de durée
\DeclareSIUnit\microliter{\micro\liter} % le modèle écrit parfois \microliter en un seul token
\DeclareSIUnit\million{millions} % dénombrement (ex. \SI{15}{\million\per\mL} spermatozoïdes)
\usepackage{qrcode}
% \degree : commande fréquemment utilisée par les modèles pour un angle ou une
% température en dehors de \SI{}{} (non fournie par les paquets chargés).
\providecommand{\degree}{^{\circ}}
% \qed (fin de démonstration), utilisé par les modèles sans amsthm
\providecommand{\qed}{\hfill$\square$}
% \barre : double barre des valeurs interdites dans un tableau de variations (tkz-tab)
\providecommand{\barre}{\ensuremath{\|}}
% environnement proof (amsthm non chargé) utilisé par les modèles
\ifdefined\proof\else\newenvironment{proof}[1][Démonstration]{\par\noindent\textit{#1.}\ }{\qed\par}\fi
\usepackage{lastpage}
\usepackage{hyperref}
\hypersetup{hidelinks}

% ============================================================
% Palette « Pop » OnBuch+ — mêmes hex que la classe P (plus_ui.dart)
% ============================================================
\definecolor{poporange}{HTML}{F59321}
\definecolor{popdark}{HTML}{E07A0C}
\definecolor{popcream}{HTML}{FFF6E8}
\definecolor{popink}{HTML}{1C1714}
\definecolor{poppurple}{HTML}{7A5AE0}
\definecolor{popgreen}{HTML}{2E9E6A}
\definecolor{poppink}{HTML}{E0457B}
\definecolor{popblue}{HTML}{2F7DE1}
\definecolor{popgold}{HTML}{C98A12}
\definecolor{popmuted}{HTML}{8A7F76}
\definecolor{popline}{HTML}{EADFD2}
\definecolor{popgray}{HTML}{8A7F76}
\colorlet{popgrey}{popgray}
% Alias tolérants (le modèle nomme parfois les couleurs différemment)
\colorlet{popwhite}{white}
\colorlet{popblack}{popink}
\colorlet{popred}{poppink}
\colorlet{popyellow}{popgold}
% Teintes claires (fonds de tableaux/encadrés) — le modèle nomme parfois
% les couleurs avec un suffixe L (ex. popblueL) sans les avoir définies.
\colorlet{poporangeL}{poporange!20}
\colorlet{popdarkL}{popdark!20}
\colorlet{poppurpleL}{poppurple!20}
\colorlet{popgreenL}{popgreen!20}
\colorlet{poppinkL}{poppink!20}
\colorlet{popblueL}{popblue!20}
\colorlet{popgoldL}{popgold!20}
\colorlet{popredL}{popred!20}
\colorlet{popgrayL}{popgray!20}
% Teintes claires pour fonds de tableaux / zones
\colorlet{popblueL}{popblue!10!white}
\colorlet{poporangeL}{poporange!14!white}
\colorlet{popgreenL}{popgreen!12!white}
\colorlet{poppurpleL}{poppurple!12!white}
\colorlet{poppinkL}{poppink!12!white}
\colorlet{popgoldL}{popgold!14!white}

\pagecolor{popcream}

% ============================================================
% Typographie
% ============================================================
\defaultfontfeatures{Ligatures=TeX}
\setmainfont{PlusJakartaSans-Medium.ttf}[Path=../../../fonts/,BoldFont=PlusJakartaSans-Bold.ttf]
% Maths en sans-serif (Fira Math) avec chiffres et lettres droites de Plus
% Jakarta, pour que les formules se fondent dans le texte.
\usepackage[math-style=ISO,bold-style=ISO]{unicode-math}
\setmathfont{FiraMath-Regular.otf}[Path=../../../fonts/]
\setmathfont{PlusJakartaSans-Medium.ttf}[Path=../../../fonts/,range={up/num,up/{Latin,latin},\mathup}]
\usepackage{icomma} % virgule décimale sans espace en mode math
% Caractères mathématiques Unicode absents des polices de texte (Plus Jakarta,
% Archivo Black) : rendus par leur équivalent mathématique, en texte comme en
% formule — sinon ils s'affichent en carré vide (ex. « Arithmétique dans ℤ »).
\usepackage{newunicodechar}
\newunicodechar{ℝ}{\ensuremath{\mathbb{R}}}
\newunicodechar{ℤ}{\ensuremath{\mathbb{Z}}}
\newunicodechar{ℂ}{\ensuremath{\mathbb{C}}}
\newunicodechar{ℕ}{\ensuremath{\mathbb{N}}}
\newunicodechar{ℚ}{\ensuremath{\mathbb{Q}}}
\newunicodechar{𝔻}{\ensuremath{\mathbb{D}}}
\newunicodechar{α}{\ensuremath{\alpha}}
\newunicodechar{β}{\ensuremath{\beta}}
\newunicodechar{γ}{\ensuremath{\gamma}}
\newunicodechar{δ}{\ensuremath{\delta}}
\newunicodechar{ε}{\ensuremath{\varepsilon}}
\newunicodechar{θ}{\ensuremath{\theta}}
\newunicodechar{λ}{\ensuremath{\lambda}}
\newunicodechar{μ}{\ensuremath{\mu}}
\newunicodechar{σ}{\ensuremath{\sigma}}
\newunicodechar{τ}{\ensuremath{\tau}}
\newunicodechar{φ}{\ensuremath{\varphi}}
\newunicodechar{ψ}{\ensuremath{\psi}}
\newunicodechar{ω}{\ensuremath{\omega}}
\newunicodechar{Σ}{\ensuremath{\Sigma}}
% majuscules produites par \MakeUppercase dans les titres (α-aminés -> Α)
\newunicodechar{Α}{\ensuremath{\alpha}}
\newunicodechar{Β}{\ensuremath{\beta}}
\newunicodechar{Γ}{\ensuremath{\Gamma}}
\newunicodechar{Δ}{\ensuremath{\Delta}}
\newunicodechar{Ε}{\ensuremath{\varepsilon}}
\newunicodechar{Θ}{\ensuremath{\Theta}}
\newunicodechar{Λ}{\ensuremath{\Lambda}}
\newunicodechar{Μ}{\ensuremath{\mu}}
\newunicodechar{Τ}{\ensuremath{\tau}}
\newunicodechar{Φ}{\ensuremath{\Phi}}
\newunicodechar{Ψ}{\ensuremath{\Psi}}
\newunicodechar{Ω}{\ensuremath{\Omega}}
\newunicodechar{↦}{\ensuremath{\mapsto}}
\newunicodechar{∈}{\ensuremath{\in}}
\newunicodechar{∉}{\ensuremath{\notin}}
\newunicodechar{∧}{\ensuremath{\wedge}}
\newunicodechar{⇔}{\ensuremath{\Leftrightarrow}}
\newunicodechar{⟺}{\ensuremath{\Longleftrightarrow}}
\newunicodechar{⇒}{\ensuremath{\Rightarrow}}
\newunicodechar{⟹}{\ensuremath{\Longrightarrow}}
\newunicodechar{∘}{\ensuremath{\circ}}
\newunicodechar{∪}{\ensuremath{\cup}}
\newunicodechar{∩}{\ensuremath{\cap}}
\newunicodechar{≡}{\ensuremath{\equiv}}
\newunicodechar{⊂}{\ensuremath{\subset}}
\newunicodechar{⊥}{\ensuremath{\perp}}
\newunicodechar{∃}{\ensuremath{\exists}}
\newunicodechar{∀}{\ensuremath{\forall}}
\newunicodechar{∛}{\ensuremath{\sqrt[3]{\,}}}
\newunicodechar{∜}{\ensuremath{\sqrt[4]{\,}}}
\newunicodechar{⃗}{\ensuremath{{}^{\to}}}
\newunicodechar{ⁿ}{\ifmmode^{n}\else\textsuperscript{\ensuremath{n}}\fi}
\newunicodechar{ˣ}{\ifmmode^{x}\else\textsuperscript{\ensuremath{x}}\fi}
\newunicodechar{ᵇ}{\ifmmode^{b}\else\textsuperscript{\ensuremath{b}}\fi}
\newunicodechar{ʳ}{\ifmmode^{r}\else\textsuperscript{\ensuremath{r}}\fi}
\newunicodechar{ᵘ}{\ifmmode^{u}\else\textsuperscript{\ensuremath{u}}\fi}
\newunicodechar{ʸ}{\ifmmode^{y}\else\textsuperscript{\ensuremath{y}}\fi}
\newunicodechar{ᵃ}{\ifmmode^{a}\else\textsuperscript{\ensuremath{a}}\fi}
\newunicodechar{ᵉ}{\ifmmode^{e}\else\textsuperscript{\ensuremath{e}}\fi}
\newunicodechar{ᵏ}{\ifmmode^{k}\else\textsuperscript{\ensuremath{k}}\fi}
\newunicodechar{ᵖ}{\ifmmode^{p}\else\textsuperscript{\ensuremath{p}}\fi}
\newunicodechar{ᴺ}{\ifmmode^{N}\else\textsuperscript{\ensuremath{N}}\fi}
\newunicodechar{ᶜ}{\ifmmode^{c}\else\textsuperscript{\ensuremath{c}}\fi}
\newunicodechar{ʰ}{\ifmmode^{h}\else\textsuperscript{\ensuremath{h}}\fi}
\newunicodechar{ᵠ}{\ifmmode^{q}\else\textsuperscript{\ensuremath{q}}\fi}
\newunicodechar{ₙ}{\ifmmode_{n}\else\textsubscript{\ensuremath{n}}\fi}
\newunicodechar{ₐ}{\ifmmode_{a}\else\textsubscript{\ensuremath{a}}\fi}
\newunicodechar{ᵢ}{\ifmmode_{i}\else\textsubscript{\ensuremath{i}}\fi}
\newunicodechar{ᵣ}{\ifmmode_{r}\else\textsubscript{\ensuremath{r}}\fi}
\newunicodechar{ₖ}{\ifmmode_{k}\else\textsubscript{\ensuremath{k}}\fi}
\newunicodechar{ᵦ}{\ifmmode_{\beta}\else\textsubscript{\ensuremath{\beta}}\fi}
\newunicodechar{ₓ}{\ifmmode_{x}\else\textsubscript{\ensuremath{x}}\fi}
\newfontfamily\popdisplay{ArchivoBlack-Regular.ttf}[Path=../../../fonts/]
\newfontfamily\poplogo{SpaceGrotesk-Bold.ttf}[Path=../../../fonts/]
\newfontfamily\popsemi{PlusJakartaSans-SemiBold.ttf}[Path=../../../fonts/]
\newfontfamily\popxbold{PlusJakartaSans-ExtraBold.ttf}[Path=../../../fonts/]
\newfontfamily\popxboldls{PlusJakartaSans-ExtraBold.ttf}[Path=../../../fonts/,LetterSpace=3.0]

\setlist[enumerate]{leftmargin=1.7em,itemsep=5pt,topsep=4pt}
\setlist[itemize]{leftmargin=1.7em,itemsep=4pt,topsep=4pt,label={\color{poporange}\textbullet}}
\setlength{\parindent}{0pt}
\setlength{\parskip}{5pt}
\renewcommand{\arraystretch}{1.25}

% pgfplots : style Pop par défaut
\pgfplotsset{
  every axis/.append style={axis line style={popink,line width=1pt},tick style={popink},
    grid style={popline,line width=0.5pt},label style={font=\small},tick label style={font=\footnotesize}},
  popcurve/.style={poporange,line width=1.8pt,no markers,smooth},
}
% Même style disponible dans un \draw TikZ simple (hors pgfplots)
\tikzset{popcurve/.style={poporange,line width=1.8pt}}
\tikzset{popgrid/.style={popline,very thin}}
\colorlet{popcurve}{poporange} % le nom du style est parfois utilisé comme couleur

% ============================================================
% Pill / squiggle
% ============================================================
\newcommand{\poppill}[3]{%
  \tikz[baseline=-0.55ex]\node[draw=popink,line width=1.2pt,rounded corners=2.6pt,%
    fill=#2,inner xsep=6.5pt,inner ysep=3pt]{%
    \popxboldls\fontsize{7}{7}\selectfont\color{#3}\MakeUppercase{#1}};%
}
\newcommand{\popsquiggle}[1]{%
  \tikz[baseline=-0.4ex]{
    \draw[#1,line width=2.6pt,line cap=round]
      plot [smooth] coordinates {(0,0.09) (0.34,0.01) (0.68,0.21) (1.02,0.01) (1.36,0.10)};
  }%
}
% Mot-clé surligné (marqueur orange sous le texte)
\newcommand{\cle}[1]{{\popxbold\color{popdark}#1}}

% ============================================================
% En-tête / pied
% ============================================================
\pagestyle{fancy}
\fancyhf{}
\renewcommand{\headrulewidth}{0pt}
\renewcommand{\footrulewidth}{0pt}
\lhead{\raisebox{-0.15ex}{\poplogo\fontsize{13}{13}\selectfont\color{popink}OnBuch\color{poporange}+}}
\rhead{\poppill{\DOCMATIERE\ \textbullet\ \DOCNIVEAU\ \textbullet\ Leçon \DOCLECON}{white}{popink}}
\lfoot{\popxbold\fontsize{7.6}{7.6}\selectfont\color{popmuted}\MakeUppercase{Cours signé OnBuch+}\ \textbullet\ {\popsemi\DOCTITRE}}
\rfoot{\tikz[baseline=-0.6ex]\node[rounded corners=8.5pt,fill=popink,minimum height=17pt,minimum width=17pt,inner xsep=5pt,inner sep=0]{\popxbold\fontsize{8}{8}\selectfont\color{popcream}\thepage\,/\,\pageref*{LastPage}};}

% ============================================================
% Titres
% ============================================================
\newcommand{\chaptitre}[2]{%
  \begin{tcolorbox}[enhanced,colback=poporange,colframe=popink,boxrule=2.6pt,arc=11pt,
      drop shadow=popink,left=15pt,right=15pt,top=12pt,bottom=11pt,before skip=3pt,after skip=18pt]
    \poppill{#2}{popink}{popcream}\par\vspace{9pt}
    {\raggedright\hyphenpenalty=10000\exhyphenpenalty=10000\popdisplay\fontsize{22}{24}\selectfont\color{popcream}\MakeUppercase{#1}\par}
  \end{tcolorbox}
}
% Section numérotée : pastille orange avec le numéro + titre Archivo
\newcounter{popsec}
\newcommand{\coursec}[1]{%
  \stepcounter{popsec}\setcounter{popsub}{0}%
  \par\vspace{16pt}\noindent
  \begin{minipage}{\linewidth}
  \tikz[baseline=-0.7ex]\node[draw=popink,line width=1.6pt,fill=poporange,rounded corners=4pt,minimum size=22pt,inner sep=0]{\popdisplay\fontsize{12}{12}\selectfont\color{popcream}\thepopsec};%
  \hspace{8pt}{\popdisplay\fontsize{15}{17}\selectfont\color{popink}\MakeUppercase{#1}}\par
  \vspace{4pt}\noindent{\color{poporange}\rule{36pt}{3.6pt}}
  \end{minipage}\par\nopagebreak\vspace{8pt}
}
\newcounter{popsub}
\newcommand{\courssub}[1]{%
  \stepcounter{popsub}%
  \par\vspace{9pt}\noindent{\popxbold\fontsize{11.5}{13}\selectfont\color{popink}{\color{poporange}\thepopsec.\thepopsub}\hspace{6pt}#1}\par\nopagebreak\vspace{4pt}
}

% ============================================================
% Boîtes pédagogiques — même recette : fond blanc, bordure épaisse colorée,
% ombre dure encre, badge plein. Titre optionnel [#1].
% ============================================================
\tcbset{popbase/.style={breakable,enhanced,colback=white,boxrule=2.3pt,arc=9pt,
  shadow={3pt}{-3pt}{0pt}{popink},left=14pt,right=14pt,top=11pt,bottom=11pt,before skip=11pt,after skip=15pt,
  fontupper=\popsemi\fontsize{10.6}{14.5}\selectfont\color{popink},
  fontlower=\popsemi\fontsize{10.6}{14.5}\selectfont\color{popink}}}
% Coche dessinée (le glyphe ✓ n'existe pas dans les polices du document)
\newcommand{\popcheck}[1]{\tikz[baseline=-0.5ex]\draw[#1,line width=1.6pt,line cap=round,line join=round](-0.13,0.0)--(-0.04,-0.09)--(0.13,0.1);}
\providecommand{\coche}{\popcheck{popgreen}}
\providecommand{\resultat}[1]{\par\smallskip\noindent\fcolorbox{popgreen}{popgreenL}{\parbox{\dimexpr\linewidth-2\fboxsep-2\fboxrule\relax}{\textbf{Résultat.} #1}}\par\smallskip}
\newcommand{\popboxhead}[3]{\poppill{#1}{#2}{white}\ifx\relax#3\relax\else\hspace{6pt}{\popxbold\fontsize{10.6}{12}\selectfont\color{popink}#3}\fi\par\vspace{7pt}}

\newtcolorbox{aretenir}[1][]{popbase,colframe=popgreen,before upper={\popboxhead{À retenir}{popgreen}{#1}}}
\newtcolorbox{exemplebox}[1][]{popbase,colframe=poporange,before upper={\popboxhead{Exemple}{poporange}{#1}}}
\newtcolorbox{definition}[1][]{popbase,colframe=popblue,before upper={\popboxhead{Définition}{popblue}{#1}}}
\newtcolorbox{propriete}[1][]{popbase,colframe=poppurple,before upper={\popboxhead{Propriété}{poppurple}{#1}}}
\newtcolorbox{methode}[1][]{popbase,colframe=popgold,before upper={\popboxhead{Méthode}{popgold}{#1}}}
\newtcolorbox{attention}[1][]{popbase,colframe=poppink,before upper={\popboxhead{Attention}{poppink}{#1}}}
\newtcolorbox{experience}[1][]{popbase,colframe=popblue,colback=popblueL,before upper={\popboxhead{Expérience}{popblue}{#1}}}
% « Le savais-tu ? » : carte violette pleine (contexte, culture, Cameroun)
\newtcolorbox{savaistu}[1][]{popbase,colback=poppurple,colframe=popink,
  fontupper=\popsemi\fontsize{10.4}{14.2}\selectfont\color{white},
  before upper={\poppill{Le savais-tu\,?}{white}{poppurple}\ifx\relax#1\relax\else\hspace{6pt}{\popxbold\color{white}#1}\fi\par\vspace{7pt}}}
% Exercice résolu : énoncé en haut, \tcblower puis la solution sur fond crème
\newcounter{popexo}\newcounter{popexr}
\newtcolorbox{exoresolu}[1][]{popbase,colframe=poporange,colbacklower=popcream,
  segmentation style={popink,line width=1.2pt,dashed},
  before upper={\stepcounter{popexr}\popboxhead{Exercice résolu \thepopexr}{poporange}{#1}},
  before lower={\poppill{Solution}{popink}{popcream}\par\vspace{6pt}}}
% Exercice (non résolu en place) — la correction est dans la partie Corrigés
\newtcolorbox{exercice}[1][]{popbase,colframe=popink,
  before upper={\stepcounter{popexo}\popboxhead{Exercice \thepopexo}{popink}{#1}}}
% Corrigé
\newtcolorbox{corrige}[1][]{popbase,colframe=popgreen,colback=popgreenL,
  before upper={\popboxhead{Corrigé}{popgreen}{#1}}}
% Objectifs / prérequis / bilan
\newtcolorbox{objectifs}{popbase,colframe=popink,colback=poporangeL,
  before upper={\popboxhead{Objectifs de la leçon}{popink}{}}}
\newtcolorbox{prerequis}{popbase,colframe=popink,colback=popblueL,
  before upper={\popboxhead{Prérequis}{popblue}{}}}
\newtcolorbox{bilan}{popbase,colframe=popink,colback=white,boxrule=2.8pt,
  before upper={\popboxhead{Fiche bilan}{poporange}{L'essentiel à connaître pour l'examen}}}
% Figure encadrée avec légende
\newtcolorbox{popfigure}[1][]{enhanced,breakable=false,colback=white,colframe=popline,boxrule=1.4pt,arc=8pt,
  left=10pt,right=10pt,top=10pt,bottom=8pt,before skip=10pt,after skip=12pt,halign=center,
  lower separated=false,halign lower=center,
  fontlower=\popsemi\fontsize{9}{11.5}\selectfont\color{popmuted},
  #1}
\newcommand{\legende}[1]{\par\vspace{6pt}{\popsemi\fontsize{9}{11.5}\selectfont\color{popmuted}#1}}

% ============================================================
% Signature OnBuch+
% ============================================================
\newcommand{\poplogoinline}[1]{{\poplogo\fontsize{#1}{#1}\selectfont\color{popink}OnBuch\color{poporange}+}}
\newcommand{\signatureonbuch}{%
  \begin{tcolorbox}[enhanced,colback=popink,colframe=popink,boxrule=2.2pt,arc=10pt,
      drop shadow=poporange,left=14pt,right=14pt,top=10pt,bottom=10pt,before skip=0pt,after skip=0pt]
    \tikz[baseline=-0.7ex]\node[circle,fill=popgreen,draw=popcream,line width=1.3pt,minimum size=22pt,inner sep=0]{\popcheck{white}};%
    \hspace{9pt}\begin{minipage}[c]{0.62\linewidth}
      {\popxbold\fontsize{12}{13}\selectfont\color{popcream}Signé\ \textbullet\ {\poplogo OnBuch\color{poporange}+}}\par\vspace{2pt}
      {\raggedright\popsemi\fontsize{8}{10.5}\selectfont\color{popline}\MakeUppercase{Équipe pédagogique OnBuch+}\\\MakeUppercase{Conforme au programme officiel MINESEC}\par}
    \end{minipage}\hfill
    {\poplogo\fontsize{22}{22}\selectfont\color{popcream}OnBuch\color{poporange}+}
  \end{tcolorbox}%
}

% ============================================================
% Page de garde
% #1 titre · #2 matière · #3 niveau · #4 module · #5 n° leçon · #6 durée ·
% #7 description / objectifs (texte ou itemize)
% ============================================================
\newcommand{\pagegarde}[7]{%
  \thispagestyle{empty}
  \newgeometry{a4paper,margin=1.7cm,top=1.6cm,bottom=1.4cm}%
  \begin{tikzpicture}[remember picture,overlay]
    \fill[poporange!18] ([xshift=-3.2cm,yshift=-3.6cm]current page.north east) circle (4.6cm);
    \fill[poppurple!14] ([xshift=2.4cm,yshift=7.6cm]current page.south west) circle (3.8cm);
  \end{tikzpicture}%
  {\poplogo\fontsize{30}{30}\selectfont\color{popink}OnBuch\color{poporange}+}\hfill
  \raisebox{0.6ex}{\poppill{Cours signé \textbullet\ Premium}{popink}{popcream}}\par
  \vspace{1pt}\popsquiggle{poppurple}\par
  \vspace{8pt}
  \begin{tcolorbox}[enhanced,colback=poporange,colframe=popink,boxrule=2.8pt,arc=13pt,
      drop shadow=popink,left=18pt,right=18pt,top=12pt,bottom=13pt,after skip=10pt]
    \poppill{#2 \textbullet\ #3}{popink}{popcream}\hspace{5pt}\poppill{#4}{white}{popink}\par\vspace{8pt}
    {\popdisplay\fontsize{14}{14}\selectfont\color{popink}LEÇON #5}\par\vspace{4pt}
    {\raggedright\hyphenpenalty=10000\exhyphenpenalty=10000\popdisplay\fontsize{25}{27}\selectfont\color{popcream}\MakeUppercase{#1}\par}
    \vspace{8pt}
    {\popxbold\fontsize{9.5}{9.5}\selectfont\color{popink}\MakeUppercase{Durée conseillée : #6}}
  \end{tcolorbox}
  \begin{tcolorbox}[enhanced,colback=white,colframe=popink,boxrule=2.2pt,arc=10pt,
      drop shadow=popink,left=16pt,right=16pt,top=10pt,bottom=10pt,after skip=8pt]
    \poppill{Dans cette leçon}{poppurple}{white}\par\vspace{5pt}
    \popsemi\fontsize{9.3}{12.6}\selectfont\color{popink}#7
  \end{tcolorbox}
  \noindent\begin{minipage}[t]{0.487\linewidth}
    \begin{tcolorbox}[enhanced,colback=popgreen,colframe=popink,boxrule=2.2pt,arc=10pt,
        drop shadow=popink,left=11pt,right=11pt,top=10pt,bottom=10pt,height=3.1cm,valign=center]
      \tcbox[colback=white,colframe=popink,boxrule=1.3pt,arc=5pt,boxsep=0pt,nobeforeafter,
        left=3pt,right=3pt,top=3pt,bottom=3pt,valign=center]{\qrcode[height=1.9cm]{https://play.google.com/store/apps/details?id=com.luvvix.onbuch&hl=fr}}%
      \hspace{8pt}\begin{minipage}[c]{0.5\linewidth}
        {\popdisplay\fontsize{11}{12.5}\selectfont\color{white}\MakeUppercase{Télécharge OnBuch}}\par\vspace{4pt}
        {\raggedright\popsemi\fontsize{8.4}{11}\selectfont\color{white}Gratuit \textbullet\ Google Play\\Tuteur IA, annales, concours, résultats\par}
      \end{minipage}
    \end{tcolorbox}
  \end{minipage}\hfill
  \begin{minipage}[t]{0.487\linewidth}
    \begin{tcolorbox}[enhanced,colback=poppurple,colframe=popink,boxrule=2.2pt,arc=10pt,
        drop shadow=popink,left=11pt,right=11pt,top=10pt,bottom=10pt,height=3.1cm,valign=center]
      \tikz[baseline=-0.7ex]\node[circle,fill=white,draw=popink,line width=1.3pt,minimum size=1.6cm,inner sep=0]{\popdisplay\fontsize{24}{24}\selectfont\color{poppurple}+};%
      \hspace{8pt}\begin{minipage}[c]{0.6\linewidth}
        {\popdisplay\fontsize{11}{12.5}\selectfont\color{white}\MakeUppercase{Passe à OnBuch+}}\par\vspace{4pt}
        {\raggedright\popsemi\fontsize{8.4}{11}\selectfont\color{white}Tous les cours signés, Tuteur IA illimité, fiches PDF — onglet OnBuch+ de l'app\par}
      \end{minipage}
    \end{tcolorbox}
  \end{minipage}
  \vspace{8pt}
  \popcontenu
  \vfill
  \signatureonbuch
  \restoregeometry
  \newpage
}

% Rangée « Ce cours contient » (page de garde)
\newcommand{\poptile}[3]{% #1 couleur #2 gros texte #3 légende
  \begin{tcolorbox}[enhanced,colback=#1,colframe=popink,boxrule=2pt,arc=9pt,shadow={2.5pt}{-2.5pt}{0pt}{popink},
      left=6pt,right=6pt,top=9pt,bottom=9pt,halign=center,valign=center,height=1.75cm,width=\linewidth,nobeforeafter]
    {\popdisplay\fontsize{12.5}{13}\selectfont\color{white}\MakeUppercase{#2}}\par\vspace{3pt}
    {\centering\popsemi\fontsize{7.8}{9.5}\selectfont\color{white}#3\par}
  \end{tcolorbox}}
\newcommand{\popcontenu}{%
  \par\noindent{\popxboldls\fontsize{7.5}{7.5}\selectfont\color{popmuted}CE COURS SIGNÉ CONTIENT}\par\vspace{7pt}
  \noindent\begin{minipage}[t]{0.235\linewidth}\poptile{popblue}{Cours}{complet, illustré, pas à pas}\end{minipage}\hfill
  \begin{minipage}[t]{0.235\linewidth}\poptile{poporange}{Méthodes}{et exercices résolus}\end{minipage}\hfill
  \begin{minipage}[t]{0.235\linewidth}\poptile{poppink}{Exercices}{type examen + corrigés}\end{minipage}\hfill
  \begin{minipage}[t]{0.235\linewidth}\poptile{popgreen}{Fiche bilan}{l'essentiel à retenir}\end{minipage}\par}

% Sommaire
\newtcolorbox{sommaire}{enhanced,colback=white,colframe=popink,boxrule=2.2pt,arc=10pt,
  drop shadow=popink,left=18pt,right=18pt,top=13pt,bottom=13pt,before skip=6pt,after skip=20pt,
  before upper={\poppill{Sommaire}{poppurple}{white}\par\vspace{9pt}\popsemi\fontsize{10.6}{14.5}\selectfont\color{popink}}}

% Signature de fin de cours — poussée tout en bas de la dernière page
\newcommand{\findecours}{%
  \par\vfill\nopagebreak
  \begin{center}\popsquiggle{poporange}\end{center}\vspace{4pt}
  \signatureonbuch
}
\AtBeginDocument{\def\llbracket{\mathopen{[\mkern-3.5mu[}}\def\rrbracket{\mathclose{]\mkern-3.5mu]}}} % intervalles d'entiers (glyphe ⟦ absent de la police)
% Glyphes absents de Fira Math : redéfinis à partir de glyphes disponibles
\AtBeginDocument{%
  \def\setminus{\mathbin{\backslash}}\let\smallsetminus\setminus
  \def\vdots{\mathord{\vbox{\baselineskip3.2pt\lineskiplimit0pt\kern4pt\hbox{.}\hbox{.}\hbox{.}}}}%
  \def\ddots{\mathinner{\mkern1mu\raise6.5pt\hbox{.}\mkern2mu\raise3.6pt\hbox{.}\mkern2mu\raise0.7pt\hbox{.}\mkern1mu}}%
  \def\triangle{\mathord{\scalebox{0.9}{$\symup{\Delta}$}}}%
  \def\nabla{\mathord{\raisebox{\depth}{\scalebox{1}[-1]{$\symup{\Delta}$}}}}%
  \def\checkmark{\popcheck{popdark}}%
  \def\star{\mathbin{\ast}}\def\bigstar{\mathord{\ast}}%
  \def\diamond{\mathbin{\scalebox{0.6}{$\Diamond$}}}%
  \def\bigcup{\mathop{\mathchoice{\raisebox{-0.3em}{\scalebox{1.7}{$\cup$}}}{\scalebox{1.25}{$\cup$}}{\cup}{\cup}}\displaylimits}%
  \def\bigcap{\mathop{\mathchoice{\raisebox{-0.3em}{\scalebox{1.7}{$\cap$}}}{\scalebox{1.25}{$\cap$}}{\cap}{\cap}}\displaylimits}%
  \def\lesseqqgtr{\mathrel{\lessgtr}}\def\bigoplus{\mathop{\oplus}\displaylimits}%
}
\providecommand{\ssi}{si et seulement si} % abréviation fréquente
\AtBeginDocument{\providecommand{\wideparen}{\overparen}} % arc de cercle

%% FIGFIX-BEGIN v5 — correctifs globaux des figures (docs/onbuch-generation/tools/figfix)
\usepackage{adjustbox}\usepackage{etoolbox}\tcbuselibrary{raster}
% toute figure TikZ/pgfplots plus large que la ligne est réduite à la largeur disponible
\BeforeBeginEnvironment{tikzpicture}{\begin{adjustbox}{max width=\linewidth}}
\AfterEndEnvironment{tikzpicture}{\end{adjustbox}}
% légendes pgfplots lisibles par-dessus les courbes
\pgfplotsset{every axis/.append style={legend style={fill=white,fill opacity=0.92,text opacity=1,draw=black!15,inner sep=3pt}}}
% étiquettes d'arêtes (syntaxe edge["..."]) avec fond blanc
\tikzset{every edge quotes/.append style={fill=white,inner sep=1.2pt}}
%% FIGFIX-END
\begin{document}

\pagegarde{La Philosophie africaine}{Philosophie}{Tle Tech}{Philosophie – classe de Terminale technique}{N01}{2 séances (fiche de progression harmonisée)}{Cette leçon explore l'histoire et les grands courants de la philosophie africaine avant d'analyser deux thématiques programmatiques majeures : le vivre-ensemble et la gouvernance. Elle vise à outiller l'élève pour mobiliser ces concepts dans l'analyse de situations sociales camerounaises actuelles.
\vspace{4pt}
\begin{multicols}{2}\small
\begin{itemize}
\item Situer chronologiquement les grandes étapes de l'émergence et de la reconnaissance de la philosophie africaine.
\item Différencier les principales tendances (ethnophilosophie, école sage, nationaliste-idéologique, critique/universaliste) et leurs représentants clés.
\item Expliquer le concept de "vivre-ensemble" à partir de ses fondements ontologiques (ubuntu, bamiléké, béti) et juridiques (Charte africaine).
\item Analyser les modèles de gouvernance traditionnels (chefferies, sociétés secrètes, consensus) et leurs apports aux défis démocratiques actuels.
\item Appliquer les notions d'unité, de consensus et de responsabilité partagée à la résolution d'un conflit local concret (ex: gestion de l'eau, succession foncière).
\item Rédiger une argumentation structurée justifiant une position éthique ou politique en s'appuyant sur des auteurs africains (Wiredu, Hountondji, Mudimbe, Oyono-Mbia).
\end{itemize}\end{multicols}}

\begin{sommaire}
\begin{multicols}{2}
\begin{enumerate}
\item 1. Genèse et historicité de la philosophie africaine
\item 2. Le "Vivre-ensemble" : fondements ontologiques et enjeux citoyens
\item 3. Gouvernance : modèles traditionnels, État moderne et défis démocratiques
\item 4. Synthèse appliquée : Articuler philosophie, vivre-ensemble et gouvernance au Cameroun
\item Activité d'intégration
\item Méthodes et exercices résolus
\item Exercices
\item Corrigés des exercices
\item Fiche bilan
\end{enumerate}
\end{multicols}
\end{sommaire}

\begin{objectifs}
\begin{itemize}
\item Situer chronologiquement les grandes étapes de l'émergence et de la reconnaissance de la philosophie africaine.
\item Différencier les principales tendances (ethnophilosophie, école sage, nationaliste-idéologique, critique/universaliste) et leurs représentants clés.
\item Expliquer le concept de "vivre-ensemble" à partir de ses fondements ontologiques (ubuntu, bamiléké, béti) et juridiques (Charte africaine).
\item Analyser les modèles de gouvernance traditionnels (chefferies, sociétés secrètes, consensus) et leurs apports aux défis démocratiques actuels.
\item Appliquer les notions d'unité, de consensus et de responsabilité partagée à la résolution d'un conflit local concret (ex: gestion de l'eau, succession foncière).
\item Rédiger une argumentation structurée justifiant une position éthique ou politique en s'appuyant sur des auteurs africains (Wiredu, Hountondji, Mudimbe, Oyono-Mbia).
\item Construire un schéma heuristique reliant ontologie africaine, éthique sociale et organisation politique.
\item Évaluer critique la pertinence de la philosophie africaine face aux enjeux de la modernité (droits de l'homme, genre, développement durable).
\end{itemize}
\end{objectifs}

\begin{prerequis}
\begin{itemize}
\item Notion de philosophie et de rationalité (programme de 1ère).
\item Repères historiques : colonisation, indépendances, construction étatique au Cameroun (Histoire 1ère/Terminale).
\item Connaissance des grands groupes culturels camerounais (Soudanais, Bantous, Sémi-bantous) et de leur organisation sociale (Géographie 1ère).
\item Maîtrise de la structure d'une dissertation philosophique (introduction, thèse/antithèse/synthèse, conclusion).
\item Notions de base en droit : Constitution, Déclaration universelle des droits de l'homme, Charte africaine des droits de l'homme et des peuples.
\end{itemize}
\end{prerequis}

\newpage

\coursec{Genèse et historicité de la philosophie africaine}

Au carrefour de la rue Mokolo à Yaoundé, un vendeur à la sauvette, un cadre de la fonction publique et un chef de quartier discutent animément de la gestion des ordures ménagères : chacun invoque tour à tour le «\cle{vivre-ensemble}», l'autorité de l'État et la solidarité du lignage pour proposer une solution. Cette scène banale révèle comment les concepts philosophiques africains structurent, sans qu'on s'en rende toujours compte, nos manières concrètes de penser la cité et le pouvoir au Cameroun.

\courssub{Problématique}
Comment la philosophie africaine, longtemps niée ou réduite à du folklore, s'est-elle constituée en discipline académique rigoureuse, capable d'éclairer les défis politiques et sociaux contemporains du Cameroun ? Nous retracerons la genèse de cette philosophie à travers quatre grands courants, de la controverse fondatrice à l'institutionnalisation universitaire camerounaise.

\courssub{Le débat inaugural : existence et nature de la philosophie africaine}

La philosophie africaine moderne naît d'un choc épistémologique. Au début du XXe siècle, l'ethnologue français \textbf{Lucien Lévy-Bruhl} soutient dans \emph{Les Fonctions mentales dans les sociétés inférieures} (1910) que l'«\cle{esprit primitif}» est \emph{pré-logique}, mystique, indifférent à la contradiction et au principe d'identité. L'Africain ne «\emph{pense}» pas, il «\emph{participe}» au monde. La philosophie, au sens grec de \emph{logos} rationnel et critique, lui serait étrangère.

\begin{definition}[Ethnophilosophie]
Étude systématique des visions du monde implicites dans les langues, mythes, institutions d'un peuple, présentée comme philosophie \textbf{collective et anonyme}. Elle ne cherche pas des auteurs individuels, mais une «\cle{philosophie implicite}» commune à toute une culture.
\end{definition}

Le missionnaire belge \textbf{Placide Tempels} retourne l'argument dans \emph{La Philosophie bantoue} (1945). Pour lui, les Bantous ont une ontologie cohérente : l'\cle{Être} est \cle{force vitale}. Toute réalité (Dieu, ancêtres, humains, nature) est force. La morale, la religion, le droit découlent de cette métaphysique unique. Tempels prouve ainsi qu'il y a \emph{une} philosophie africaine, mais il la fige dans une essence immuable et collective.

\begin{attention}[Négritude vs Philosophie africaine]
Ne pas confondre la \textbf{Négritude} (mouvement littéraire et politique francophone des années 1930-1960, porté par Senghor, Césaire, Damas) et la \textbf{Philosophie africaine} (discipline académique argumentée, soumise aux normes de la rationalité universelle). Senghor fait le pont : il «\emph{philosophise}» la Négritude (l'émotion est nègre, la raison hellène), mais l'ethnophilosophie de Tempels n'est pas l'œuvre de Senghor.
\end{attention}

\courssub{La période fondatrice : l'ethnophilosophie et la quête d'identité (1945-1970)}

Dans le sillage de Tempels, l'abbé rwandais \textbf{Alexis Kagamé} (\emph{La Philosophie bantu-rwandaise de l'être}, 1956) analyse les catégories grammaticales du kinyarwanda pour en déduire une métaphysique : l'Être se divise en \emph{ntu} (force vitale), \emph{hantu} (lieu), \emph{kantu} (chose), \emph{untu} (qualité). C'est une \cle{philosophie de l'être} déduite de la structure linguistique.

\begin{aretenir}[Caractères de l'ethnophilosophie]
\begin{itemize}
    \item \textbf{Anonymat :} pas d'auteur nommé, «\emph{le peuple}» pense.
    \item \textbf{Unanimisme :} consensus supposé sans contradiction interne.
    \item \textbf{Ahistoricité :} structures présentées comme éternelles, hors du temps.
    \item \textbf{Apologétique :} vise à prouver l'égalité de dignité face au colonisateur.
\end{itemize}
\end{aretenir}

Cette démarche a eu le mérite historique de \textbf{réhabiliter la pensée africaine} face au déni colonial. Elle a fourni un socle identitaire aux mouvements d'indépendance. Mais elle fige la tradition et empêche toute critique interne.

\courssub{La rupture épistémologique : philosophie des Sages et école critique professionnelle (1970-1990)}

Dans les années 1970, deux réponses distinctes surgissent contre l'ethnophilosophie. Il ne faut pas les confondre, bien qu'elles partagent l'exigence de rigueur.

\begin{itemize}
    \item \textbf{La philosophie des Sages (Sage Philosophy) : Henry Odera Oruka} (Kenya), \emph{Trends in Contemporary African Philosophy} (1990). Il 인터뷰 des sages traditionnels (comme \emph{Paul Mbuya Akoko}) pour extraire leurs pensées \emph{individuelles}, critiques et argumentées sur l'existence, Dieu, la connaissance. Il prouve que la tradition porte des \cle{philosophes-sages} identifiables, non un bloc anonyme.
    \item \textbf{L'école critique professionnelle (dite «\cle{école des Sages}» par abus de langage) :} Philosophes formés aux standards occidentaux (logique analytique, phénoménologie, marxisme) qui refusent toute spécificité essentialiste.
    \begin{itemize}
        \item \textbf{Paulin Hountondji} (Bénin), \emph{Sur la philosophie africaine} (1976) : distingue \emph{philosophie africaine} (ensemble des œuvres produites par des Africains) et \emph{ethnophilosophie} (discours sur la philosophie traditionnelle). Il exige une \cle{philosophie écrite}, \cle{signée}, \cle{argumentée}, soumise à la critique universelle. «\emph{Il n'y a de philosophie que dans le dialogue critique}».
        \item \textbf{Kwasi Wiredu} (Ghana), \emph{Philosophy and an African Culture} (1980) : préconise la \cle{décolonisation conceptuelle}. Il faut purger les langues africaines des concepts importés (ex: «\emph{âme}» chrétien) pour retrouver des concepts opératoires (ex: le \emph{sunsum} akan comme principe vital, non comme substance immortelle). Il défend un \cle{universalisme} : la logique est une, la vérité est une.
        \item \textbf{V.Y. Mudimbe} (RDC/USA), \emph{L'Invention de l'Afrique} (1988) : montre comment le discours occidental (missionnaires, ethnologues, coloniaux) a «\emph{inventé}» l'Afrique comme altérité radicale. L'ethnophilosophie elle-même est un produit de cette invention coloniale.
    \end{itemize}
\end{itemize}

\begin{methode}[Comparer trois courants philosophiques : tableau à double entrée]
Pour opposer rigoureusement l'ethnophilosophie, la philosophie des Sages et l'école critique, construisez un tableau à triple entrée.
\begin{enumerate}
    \item \textbf{Lignes :} Critères d'analyse (Rapport à la tradition, Visée universaliste/particulariste, Méthode, Statut de l'individu, Projet politique).
    \item \textbf{Colonnes :} Les trois courants comparés.
    \item \textbf{Remplissage :} Citez un auteur et une œuvre de référence par case.
    \item \textbf{Synthèse :} Dégagez la double rupture : passage du collectif à l'individuel (Oruka), puis de l'implicite à l'explicite critique (Hountondji/Wiredu).
\end{enumerate}
\end{methode}

\begin{exemplebox}[Application du tableau : Les trois visages de la rupture]
\begin{center}
\begin{tabularx}{\linewidth}{|l|X|X|X|}\hline
\textbf{Critère} & \textbf{Ethnophilosophie (Tempels, Kagamé)} & \textbf{Phil. des Sages (Oruka)} & \textbf{École critique (Hountondji, Wiredu)} \\ \hline
Rapport à la tradition & Source unique, sacralisée, immuable & Matériau vivant, porté par des experts identifiés & Matériau à critiquer, historiciser, réinventer \\ \hline
Visée & Particulariste (prouver l'altérité) & Particulariste assumée (sagesse locale) & Universaliste (dialogue avec la raison humaine) \\ \hline
Méthode & Inductive (langue $\to$ métaphysique), descriptive & Empirique (entretiens), herméneutique & Analytique, argumentative, normative \\ \hline
Sujet pensant & Collectif anonyme («\emph{le Bantu}») & Individuel nommé (le Sage), oral mais rigoureux & Individuel, nommé, responsable de ses thèses écrites \\ \hline
Projet politique & Réhabilitation identitaire (fierté) & Valorisation du patrimoine oral & Libération intellectuelle (autonomie critique) \\ \hline
\end{tabularx}
\end{center}
\legende{Tableau comparatif des trois courants majeurs de la période 1945-1990.}
\end{exemplebox}

\courssub{Le courant nationaliste-idéologique : philosophie de la libération et projet de société}

Parallèlement aux débats académiques, des chefs d'État intellectuels font de la philosophie le \cle{fondement idéologique} de la nouvelle nation. La philosophie quitte l'université pour la place publique.

\begin{itemize}
    \item \textbf{Kwame Nkrumah} (Ghana), \emph{Consciencisme} (1964) : synthétise trois courants (traditionnel, islamo-arabe, euro-chrétien) en une idéologie révolutionnaire : le \cle{consciencisme}. La philosophie est arme de la \cle{décolonisation mentale} et base du panafricanisme politique.
    \item \textbf{Julius Nyerere} (Tanzanie), \emph{Ujamaa} (1967) : le \cle{socialisme africain} ne vient pas de Marx mais de la structure familiale traditionnelle (\emph{ujamaa} = parenté). La philosophie justifie la villagisation et l'autosuffisance.
    \item \textbf{Léopold Sédar Senghor} (Sénégal), \emph{Liberté I-V} : la \cle{Négritude} devient philosophie de la civilisation de l'Universel. L'émotion nègre et la raison grecque se complètent pour bâtir une société communautaire moderne.
\end{itemize}

\begin{attention}[Piège : essentialisme politique]
Ces philosophies d'État risquent de figer la tradition en «\emph{coutume}» officielle pour légitimer le parti unique. La critique de Hountondji s'applique aussi ici : toute philosophie qui sert d'idéologie d'État sans espace de dissentiment cesse d'être philosophie pour devenir \emph{propagande}.
\end{attention}

\courssub{L'émergence du courant critique/universaliste contemporain : déconstruction, historicité, dialogue}

Depuis les années 1990, une nouvelle génération dépasse l'alternative «\emph{tradition vs modernité}». Elle assume l'hybridité, l'historicité et la pluralité.

\begin{itemize}
    \item \textbf{Achille Mbembe} (Cameroun/France/USA), \emph{De la postcolonie} (2000), \emph{Critique de la raison nègre} (2013) : analyse le pouvoir postcolonial comme \cle{convivialité} et \cle{simulacre}. Il déconstruit la «\emph{raison nègre}» comme invention coloniale et propose une \cle{planétarité} africaine : penser le monde depuis l'Afrique, non pour l'Afrique seulement.
    \item \textbf{Souleymane Bachir Diagne} (Sénégal/USA), \emph{En quête d'Afrique(s)} (2011) : la philosophie africaine est une \cle{philosophie en situation}. Elle fait de la \cle{traduction} (au sens de Paul Ricœur) son opérateur central : traduire les concepts traditionnels dans le langage de la modernité et vice-versa.
    \item \textbf{Fabien Eboussi Boulaga} (Cameroun), \emph{Christianisme sans fétiche} (1981), \emph{La Crise du Muntu} (1984) : jésuite devenu philosophe laïc, il opère une \cle{théologie de la libération} philosophique. Il refuse le «\emph{Muntu}» (l'homme essentialisé de Tempels) pour l'\cle{historique} : l'Africain est un sujet qui \emph{fait} son histoire dans la contradiction et le travail.
\end{itemize}

\begin{savaistu}[Fabien Eboussi Boulaga (1934-2018) : le pionnier camerounais]
Né à Bafia, jésuite, docteur en philosophie à la Sorbonne, il publie \emph{Christianisme sans fétiche} (1981). Il y démontre que le christianisme importé fonctionne comme un \cle{fétiche} (objet de pouvoir magique) s'il ne passe pas par la critique rationnelle et l'inculturation vraie. Il fonde une \textbf{philosophie africaine de la praxis} : l'homme africain n'est pas une «\emph{essence}» (Muntu) mais un \emph{projet historique}. Il a enseigné à l'Université de Yaoundé I et marqué des générations d'étudiants par son exigence conceptuelle et son engagement éthique.
\end{savaistu}

\begin{exemplebox}[Chiffres clés : Vitalité éditoriale actuelle]
En 2023, plus d'une \textbf{quinzaine de revues spécialisées} en philosophie africaine existent sur le continent (ex: \emph{Philosophia Africana}, \emph{Journal of African Philosophy}, \emph{Pensée et Développement}, \emph{Revue Internationale de Philosophie Africaine}). Cela prouve l'institutionnalisation réussie du courant critique : la philosophie africaine produit désormais ses propres normes de validation scientifique.
\end{exemplebox}

\courssub{L'institutionnalisation au Cameroun : de l'UY\_1 à l'ACAP}

Le Cameroun est un laboratoire privilégié de cette histoire.

\begin{itemize}
    \item \textbf{Années 1970} : Structuration du \textbf{Département de Philosophie} à l'Université de Yaoundé (future UY\_1). Premiers enseignants : Eboussi Boulaga, Jean-Marc Ela, Pierre Mvogo, puis plus tard Maurice Tchuente, Jean-Godefroy Bidima.
    \item \textbf{Années 1980-1990} : Émergence de l'\textbf{École de Yaoundé} (anthropologie philosophique, éthique sociale, philosophie politique). Publication de la revue \emph{Pensée et Développement} (fondée 1973, relancée par le département).
    \item \textbf{1992} : Création de l'\textbf{ACAP} (Association Camerounaise de Philosophie). Elle organise des colloques biannuels, publie les \emph{Actes de l'ACAP}, structure la recherche (équipes : éthique, épistémologie, politique, interculturalité).
    \item \textbf{Aujourd'hui} : Le département forme des docteurs qui intègrent les universités d'Afrique centrale et la diaspora. La philosophie technique (appliquée à la gouvernance, l'environnement, la santé) y trouve un terrain d'application concret.
\end{itemize}

\begin{exoresolu}[Analyse d'une citation : identifier le courant]
\textbf{Énoncé :} «\emph{La philosophie africaine ne saurait être une ethnophilosophie, c'est-à-dire une pensée anonyme, collective, figée dans la tradition. Elle doit être une philosophie critique, écrite, signée, capable de dialogue avec les autres philosophies du monde.}» (Inspiré de Hountondji, \emph{Sur la philosophie africaine}, 1976).

\textbf{Question :} À quel courant appartient cet extrait ? Justifiez en mobilisant trois critères du tableau comparatif.

\tcblower
\textbf{Solution détaillée :}
\begin{enumerate}
    \item \textbf{Identification :} Courant de l'\textbf{école critique professionnelle (Hountondji, Wiredu)}.
    \item \textbf{Justification par trois critères :}
    \begin{itemize}
        \item \textbf{Statut de l'individu :} «\emph{écrite, signée}» $\to$ sujet pensant individuel, nommé, responsable (vs anonyme collectif).
        \item \textbf{Méthode :} «\emph{critique}» $\to$ démarche argumentative, normative, soumise à la raison universelle (vs descriptive, inductive sur la langue).
        \item \textbf{Visée :} «\emph{dialogue avec les autres philosophies du monde}» $\to$ universalisme, ouverture interculturelle (vs particularisme identitaire, clôture sur la tradition).
    \end{itemize}
    \item \textbf{Conclusion :} Ce texte est le manifeste fondateur de la rupture épistémologique des années 1970 contre l'ethnophilosophie.
\end{enumerate}
\end{exoresolu}

\begin{popfigure}
\begin{tikzpicture}[
    timeline/.style={very thick, popblue},
    period/.style={fill=popblue!20, draw=popblue, rounded corners=3pt, inner sep=4pt, font=\small\bfseries, align=center},
    author/.style={fill=poporange!20, draw=poporange, rounded corners=2pt, inner sep=2pt, font=\tiny, align=center},
    event/.style={fill=popgreen!20, draw=popgreen, rounded corners=2pt, inner sep=2pt, font=\tiny, align=center},
    arrow/.style={-Stealth, popmuted, thick}
]
% Axe principal
\draw[timeline] (0,0) -- (16,0) node[right, font=\small\bfseries] {Années};
% Graduations décennales
\foreach \x/\year in {0/1945, 1.6/1950, 3.2/1960, 4.8/1970, 6.4/1980, 8/1990, 9.6/2000, 11.2/2010, 12.8/2020, 14.4/2025} {
    \draw (\x, -0.1) -- (\x, 0.1) node[below, font=\tiny] {\year};
}

% COURANT 1 : Ethnophilosophie (1945-1970)
\node[period, align=center] (ethno) at (2.4, 2.0){Ethnophilosophie\\ (Fondation identitaire)};
\draw[arrow] (ethno.south) -- (1.6, 0.2);
\node[author, align=center] at (0.8, 1.3){Tempels:\\ \emph{Phil. bantoue} (1945)};
\node[author, align=center] at (2.0, 2.8){Kagamé:\\ \emph{Être bantu} (1956)};
\node[author, align=center] at (3.2, 1.3){Senghor:\\ Négritude (ann. 30-60)};

% COURANT 2 : Nationaliste-idéologique (1957-1975)
\node[period, align=center] (nat) at (4.8, -2.0){Nationaliste-\\idéologique\\ (Libération/État)};
\draw[arrow] (nat.north) -- (4.8, -0.2);
\node[author, align=center] at (3.5, -1.3){Nkrumah:\\ Consciencisme (1964)};
\node[author, align=center] at (4.8, -2.8){Nyerere:\\ Ujamaa (1967)};
\node[event, align=center] at (6.0, -1.3){Indépendances\\ (1960)};

% COURANT 3 : Philosophie des Sages / École critique (1970-1990)
\node[period, align=center] (crit) at (7.2, 2.0){Rupture épistémologique\\ Sages (Oruka) \& Critique\\ (Hountondji, Wiredu)};
\draw[arrow] (crit.south) -- (6.4, 0.2);
\node[author, align=center] at (5.6, 1.3){Oruka:\\ \emph{Sage Phil.} (1990)};
\node[author, align=center] at (7.2, 2.8){Hountondji:\\ \emph{Sur la phil.} (1976)};
\node[author, align=center] at (8.8, 1.3){Wiredu:\\ \emph{Decol. Concept.} (1980)};
\node[author, align=center] at (8.8, 2.8){Mudimbe:\\ \emph{Inv. Afrique} (1988)};
\node[event, align=center] at (6.4, 3.2){Création Dépt Philo\\ UY\_1 (ann. 70)};

% COURANT 4 : Contemporain universaliste (1990-...)
\node[period, align=center] (contemp) at (10.4, -2.0){Critique/\\Universaliste\\ contemporain\\(Historicité/Dialogue)};
\draw[arrow] (contemp.north) -- (9.6, -0.2);
\node[author, align=center] at (9.0, -1.3){Eboussi Boulaga:\\ \emph{Chr. sans fétiche} (1981)};
\node[author, align=center] at (10.4, -2.8){Mbembe:\\ \emph{Postcolonie} (2000)};
\node[author, align=center] at (11.8, -1.3){Diagne:\\ \emph{En quête d'Afr.} (2011)};
\node[event] at (9.6, -3.2) {Conf. de Baule (1990)};
\node[event, align=center] at (11.2, -3.2){Printemps africain\\ (1990s)};
\node[event, align=center] at (12.8, -3.2){ACAP (1992)\\ Revues spécialisées};

% Flèches de continuité/rupture
\draw[poppurple, dashed, thick, -Stealth] (3.2, 1.7) to[bend left=30] node[midway, above, font=\tiny, poppurple] {Rupture : Subjectivité} (5.6, 1.7);
\draw[poppurple, dashed, thick, -Stealth] (7.2, -1.7) to[bend right=30] node[midway, below, font=\tiny, poppurple] {Continuité critique} (9.0, -1.7);

% Titre
\node[font=\bfseries\large, popdark] at (8, 3.8) {Frise chronologique : Genèse et historicité de la philosophie africaine (1945-2025)};
\end{tikzpicture}
\legende{Frise synthétique positionnant les 4 courants majeurs, auteurs clés, événements politiques et institutionnalisation camerounaise. Les flèches violettes pointillées indiquent les ruptures (ethnophilosophie $\to$ Sages/Critique) et continuités (critique $\to$ contemporain).}
\end{popfigure}

\begin{aretenir}[Synthèse de la section : 4 temps, 4 visages de la philosophie africaine]
\begin{enumerate}
    \item \textbf{1945-1970 : L'Ethnophilosophie} (Tempels, Kagamé). \emph{Projet :} Prouver l'existence d'une pensée africaine cohérente. \emph{Limite :} Essentialisme, anonymat, ahistoricité.
    \item \textbf{1960-1980 : Le Nationalisme idéologique} (Nkrumah, Nyerere, Senghor). \emph{Projet :} Penser la nation, légitimer l'État postcolonial. \emph{Risque :} Instrumentalisation politique, pensée unique.
    \item \textbf{1970-1990 : La Rupture épistémologique} (Oruka, Hountondji, Wiredu, Mudimbe). \emph{Double visage :} Philosophie des Sages (subjectivité orale identifiée) et École critique professionnelle (rigueur écrite universaliste). \emph{Apport :} Professionnalisation académique, fin de l'unanimisme.
    \item \textbf{1990-\ldots : Le Courant contemporain} (Mbembe, Diagne, Eboussi Boulaga). \emph{Projet :} Historicité, hybridité, traduction, planétarité. \emph{Ancrage camerounais :} UY\_1, ACAP, \emph{Pensée et Développement}.
\end{enumerate}
\end{aretenir}

\coursec{Le « Vivre-ensemble » : fondements ontologiques et enjeux citoyens}

Dans la section précédente, nous avons montré que la philosophie africaine, après avoir dû conquérir son droit à l'existence face aux préjugés colonialistes, s'est structurée en tendances (ethnophilosophie, philosophie de la sagesse, philosophie nationaliste-idéologique, philosophie critique) et s'ouvre aujourd'hui à de nouvelles thématiques. Parmi elles, deux questions dominent le débat public camerounais : comment vivre ensemble dans une société plurielle, et comment bien gouverner ? Ces deux questions ne sont pas étrangères à la philosophie : elles en sont le cœur, car elles interrogent la nature de la personne humaine et le fondement du lien social. C'est ce que nous examinons maintenant.

\courssub{L'ontologie relationnelle : « Je suis parce que nous sommes »}

\textbf{Intuition de départ.} Observez un enfant dans un village camerounais : avant même de savoir parler, il est déjà « fils de », « neveu de », « petit-fils de ». Son nom lui-même est une mémoire familiale. Son identité ne précède pas ses relations : elle en est constituée. C'est cette intuition que la pensée africaine classique élève au rang de thèse philosophique.

\begin{definition}[Ubuntu / Buntou]
L'\cle{Ubuntu} (langues nguni, Afrique australe) ou \cle{Buntou} (racine bantu : \emph{bu-}, préfixe de classe abstraite, et \emph{-ntu}, l'humain) est une philosophie éthique centrée sur l'\cle{humanité partagée}, l'\cle{interdépendance}, la \cle{compassion} et la \cle{réconciliation}. Sa formule canonique est : « Je suis parce que nous sommes » (\emph{Umuntu ngumuntu ngabantu} : « une personne est une personne à travers les autres personnes »). Elle a inspiré la Commission Vérité et Réconciliation en Afrique du Sud (instaurée en 1995, travaux à partir de 1996), présidée par Desmond Tutu, qui a préféré la réconciliation à la vengeance après l'apartheid.
\end{definition}

\textbf{L'énoncé rigoureux.} En philosophie occidentale moderne, le sujet se fonde sur lui-même : Descartes, dans le \emph{Discours de la méthode} (1637), affirme « \emph{Cogito, ergo sum} » — « Je pense, donc je suis ». Le « je » est premier, solitaire, substance pensante qui n'a besoin de rien d'autre pour exister. La communauté vient ensuite, comme un contrat entre individus déjà constitués. La pensée africaine classique inverse le rapport : c'est la relation qui fonde la personne. Le philosophe ghanéen Kwasi \cle{Wiredu} (\emph{Person and Community in African Traditional Thought}) et le théologien kényan John \cle{Mbiti} (« Je suis parce que nous sommes, et parce que nous sommes, je suis ») montrent que la \cle{personne} africaine est un \textbf{nœud de relations} : relations familiales et claniques (lignage), relations spirituelles (ancêtres, Dieu), relations cosmiques (la nature habitée). Devenir une personne est un \emph{processus} : on naît humain biologiquement, mais on devient \emph{personne} socialement, par les rites d'initiation, les responsabilités assumées et la reconnaissance du groupe.

\begin{popfigure}
\begin{center}
\begin{tikzpicture}[font=\small, every node/.style={align=center}]
\node[draw=popdark, fill=popblueL, rounded corners, thick, minimum width=3.4cm, minimum height=1.1cm, font=\bfseries, align=center] (P) at (0,0){Personne\\africaine};
\node[draw=popgreen, fill=popgreenL, rounded corners, thick, minimum width=3.2cm, align=center] (F) at (-4.6,2.4){1. Famille / Clan\\ \emph{Solidarité}};
\node[draw=poporange, fill=poporangeL, rounded corners, thick, minimum width=3.2cm, align=center] (V) at (4.6,2.4){2. Village / Cité\\ \emph{Participation}};
\node[draw=poppurple, fill=poppurpleL, rounded corners, thick, minimum width=3.2cm, align=center] (A) at (-4.6,-2.4){3. Ancêtres / Esprits\\ \emph{Continuité}};
\node[draw=popgold, fill=popgoldL, rounded corners, thick, minimum width=3.2cm, align=center] (N) at (4.6,-2.4){4. Nature / Cosmos\\ \emph{Harmonie}};
\draw[<->, >=Stealth, very thick, popgreen] (P) -- (F);
\draw[<->, >=Stealth, very thick, poporange] (P) -- (V);
\draw[<->, >=Stealth, very thick, poppurple] (P) -- (A);
\draw[<->, >=Stealth, very thick, popgold] (P) -- (N);
\draw[<->, >=Stealth, dashed, popmuted] (F) -- (V);
\draw[<->, >=Stealth, dashed, popmuted] (A) -- (N);
\end{tikzpicture}
\end{center}
\legende{Schéma heuristique de l'ontologie relationnelle : la personne africaine est constituée par quatre cercles de relations interdépendantes (flèches bidirectionnelles : chaque pôle donne à la personne son identité et reçoit d'elle des devoirs).}
\end{popfigure}

\begin{propriete}[Primauté du devoir sur le droit]
Dans la pensée africaine classique, le droit est \cle{subsidié au devoir} : on possède des droits \emph{parce qu'on} remplit ses devoirs envers le groupe (Wiredu, \emph{Person and Community in African Traditional Thought}). L'enfant qui honore ses parents, le notable qui nourrit la communauté, le chef qui protège le village acquièrent ainsi une dignité et des prérogatives reconnues. La personne pleinement accomplie — l'« ancêtre » en puissance — est celle qui a le plus donné au lien communautaire.
\end{propriete}

\begin{attention}[Piège de dissertation]
Ne présentez jamais l'opposition Ubuntu/Cogito comme un conflit entre une pensée « supérieure » et une pensée « inférieure ». Le cogito cartésien a fondé l'autonomie du sujet et les droits de l'homme ; l'Ubuntu rappelle que l'autonomie sans lien social devient isolement. La bonne démarche est la \emph{mise en tension} : l'identité personnelle suppose-t-elle d'abord la conscience de soi ou la reconnaissance d'autrui ?
\end{attention}

\courssub{Les fondements camerounais du vivre-ensemble}

L'ontologie relationnelle n'est pas une abstraction : elle s'incarne dans des institutions concrètes, différentes selon les aires culturelles du Cameroun, mais animées par la même logique de solidarité.

\begin{exemplebox}[Quatre institutions camerounaises du lien social]
\begin{itemize}
\item Le \cle{Njangi} (tontine pratiquée notamment chez les Bamiléké, les Bassa et bien d'autres groupes) : association rotative d'épargne et de crédit où chacun cotise et reçoit à tour de rôle. Au-delà de l'économie, le njangi est une école de confiance, de parole donnée et d'entraide : manquer à son tour, c'est se déshonorer devant le groupe.
\item Le \cle{Mvett} (épopée chantée des Fang et des Beti) : récit héroïque accompagné de la harpe-cithare, qui transmet l'éthique guerrière, la généalogie et les valeurs du lignage. Le mvet est une « université orale » : il éduque au courage, à la loyauté et à la mémoire commune.
\item Le \cle{Dina} (coutume consensuelle des Peul et des Gbaya) : ensemble de règles discutées et acceptées par la communauté, révisables par l'assemblée. Le dina montre que la tradition africaine n'est pas figée : elle est un droit vivant, produit du consentement.
\item Le \cle{Ngondo} (célébration unificatrice des Sawa, au bord du Wouri) : rassemblement annuel des peuples côtiers autour des rites aquatiques, du culte des ancêtres et de la fête commune. Le Ngondo dépasse les clans : il fabrique une identité collective partagée.
\end{itemize}
\end{exemplebox}

\textbf{Ce que ces institutions prouvent.} Leur diversité même est un argument : le vivre-ensemble camerounais ne repose pas sur une culture unique, mais sur une \emph{grammaire commune} — la solidarité, la parole, la médiation des anciens, la fête partagée — déclinée en langues différentes. C'est exactement la définition de l'unité dans la diversité.

\begin{savaistu}[La palabre, patrimoine de l'humanité]
Le concept de \cle{palabre} (du portugais \emph{palavra}, « parole ») désigne la pratique africaine de discussion publique pour prévenir ou résoudre les conflits : on parle longuement, sous l'arbre à palabres, jusqu'à ce que la parole réconcilie. L'UNESCO a inscrit des pratiques de ce type au patrimoine culturel immatériel de l'humanité, reconnaissant la palabre comme un instrument universel de médiation.
\end{savaistu}

\courssub{Le vivre-ensemble comme valeur juridique : droits et devoirs corrélatifs}

\textbf{Du village à la République.} La question philosophique devient politique : comment transposer l'ontologie relationnelle à l'échelle d'un État de plus de 250 groupes ethniques ? Le Cameroun a choisi d'inscrire le vivre-ensemble dans sa norme suprême.

\begin{aretenir}[Le vivre-ensemble dans la Constitution]
Le Préambule de la Constitution du Cameroun du 18 janvier 1996 (révisée en 2008) affirme l'attachement du peuple camerounais aux libertés fondamentales et au \cle{vivre-ensemble}. Le vivre-ensemble n'est donc pas un vœu moral facultatif : c'est une \textbf{valeur constitutionnelle} qui engage l'État et chaque citoyen.
\end{aretenir}

\textbf{La dimension panafricaine.} La Charte africaine des droits de l'homme et des peuples (adoptée à Nairobi en 1981, entrée en vigueur en 1986) est la traduction juridique continentale de l'ontologie relationnelle. Son originalité, par rapport aux textes occidentaux, est d'énoncer des \textbf{devoirs} à côté des droits :

\begin{itemize}
\item \textbf{Article 27} : tout individu a des devoirs envers sa famille, la société, l'État et la communauté internationale ; les droits s'exercent dans le respect des droits d'autrui, de la sécurité collective, de la morale et de l'intérêt commun.
\item \textbf{Article 28} : tout individu a le devoir de respecter son prochain sans discrimination et de maintenir avec lui des relations de tolérance, de dialogue et de concertation.
\item \textbf{Article 29} : devoirs de préserver la famille, de respecter ses parents, de servir la communauté nationale, de préserver l'indépendance et l'intégrité territoriale, de contribuer à la défense de la patrie et de payer les impôts.
\end{itemize}

On retrouve ici, juridiquement formulée, la propriété de Wiredu : les droits et les devoirs sont \cle{corrélatifs}. Être citoyen camerounais, c'est à la fois exiger ses droits (éducation, sécurité, liberté d'expression) et assumer ses devoirs (respect d'autrui, contribution à l'impôt, défense de l'intérêt commun).

\begin{exemplebox}[L'urgence éducative en chiffres]
Le Cameroun compte \textbf{plus de 250 groupes ethniques}, \textbf{2 langues officielles} (français et anglais), et une population très jeune : \textbf{plus de la moitié des Camerounais ont moins de 20 ans} et la grande majorité a moins de 35 ans (projections de l'INS). Cette jeunesse massive, hyperconnectée et exposée aux discours de division, fait de l'éducation au vivre-ensemble une urgence nationale : c'est à l'école que se joue la cohésion de demain.
\end{exemplebox}

\courssub{Les menaces contemporaines sur le vivre-ensemble}

\textbf{Diagnostic philosophique.} Si le vivre-ensemble est un pacte, il peut être rompu. Quatre menaces majeures pèsent aujourd'hui sur le pacte social camerounais :

\begin{enumerate}
\item \textbf{Le tribalisme politique} : lorsque l'appartenance ethnique devient un critère d'accès aux postes, aux marchés publics ou aux votes, la citoyenneté commune se fragmente en clientèles. Le groupe prime sur la personne — c'est la \emph{perversion} de l'ontologie relationnelle, qui visait l'ouverture et non l'enfermement.
\item \textbf{Les discours de haine sur les réseaux sociaux} : la parole, qui dans la palabre réconciliait, devient arme de stigmatisation. L'anonymat numérique dissout la responsabilité que supposait la parole publique face au groupe.
\item \textbf{Les crises sécuritaires} : la crise anglophone (régions du Nord-Ouest et du Sud-Ouest) et les exactions de Boko Haram (Extrême-Nord) montrent ce que coûte la rupture du dialogue : des milliers de déplacés, des écoles fermées, des vies brisées. Elles rappellent que le vivre-ensemble exige l'écoute des griefs avant qu'ils ne deviennent des armes.
\item \textbf{La corruption} : détourner le bien commun, c'est voler la communauté — l'exact contraire du njangi, où l'argent collecté sert à tous. La corruption est une rupture du pacte social car elle détruit la confiance, condition de tout lien.
\end{enumerate}

\begin{attention}[Unité n'est pas uniformité]
Le « vivre-ensemble » n'est pas la pensée unique ni l'effacement des différences culturelles : c'est la \textbf{gestion pacifique de la différence}. Confondre unité et uniformité est une erreur fréquente en dissertation : un État qui imposerait une seule langue, une seule culture, une seule opinion ne réaliserait pas le vivre-ensemble, il le détruirait par l'étouffement. L'unité camerounaise est une \emph{unité composée}.
\end{attention}

\courssub{La pédagogie du vivre-ensemble : construire la cohésion}

Le vivre-ensemble ne se décrète pas : il s'éduque et s'organise. Le Cameroun dispose de plusieurs instruments :

\begin{itemize}
\item \textbf{L'Éducation à la citoyenneté et à la morale (ECM)} : discipline scolaire qui forme au respect d'autrui, aux droits et devoirs du citoyen, à la culture de la paix.
\item \textbf{Le bilinguisme effectif} : maîtriser le français \emph{et} l'anglais, ce n'est pas seulement un atout professionnel ; c'est pouvoir parler à l'autre moitié du pays, donc reconnaître l'autre comme semblable.
\item \textbf{La décentralisation} : en rapprochant le pouvoir des citoyens (régions, communes), elle permet la participation locale — l'équivalent moderne du cercle du village.
\item \textbf{La Commission nationale pour la promotion du bilinguisme et du multiculturalisme (CNPBM)} : institution créée en 2017 pour veiller à l'équilibre linguistique et culturel et promouvoir la cohésion nationale.
\end{itemize}

\begin{methode}[Analyser un conflit local avec les outils de la leçon]
\begin{enumerate}
\item Identifier les \textbf{acteurs} et leurs appartenances (familiales, ethniques, religieuses, politiques).
\item Déterminer la \textbf{norme violée} : s'agit-il de la coutume (dina, droit lignager), de la loi étatique, ou des deux ?
\item Repérer le \textbf{mécanisme traditionnel de résolution} disponible : palabre, médiation des anciens, chef traditionnel, rite de réconciliation.
\item Articuler avec la \textbf{loi étatique} : que dit la Constitution ? Quelle juridiction est compétente ? Comment concilier justice coutumière et justice républicaine sans contradiction ?
\end{enumerate}
\end{methode}

\begin{exoresolu}[Conflit foncier entre deux villages]
\textbf{Énoncé.} Deux villages voisins de la région de l'Ouest camerounais se disputent un terrain de 12 hectares. Des jeunes du village A y ont planté des arbres ; les notables du village B brandissent un titre foncier obtenu auprès de l'administration. Des menaces circulent sur WhatsApp. Appliquez la méthode d'analyse en quatre étapes et proposez une issue conforme au vivre-ensemble.
\tcblower
\textbf{Solution détaillée.}
\begin{enumerate}
\item \textbf{Acteurs et appartenances} : les jeunes planteurs du village A (génération montante, usage coutumier de la terre), les notables du village B (détenteurs d'un titre administratif), les chefs traditionnels des deux lignages, l'administration (sous-préfet), et les relais numériques anonymes qui attisent le conflit.
\item \textbf{Norme violée} : double référentiel. Pour la coutume, la terre appartient au lignage qui l'a défrichée et transmise — la mise en valeur par A crée un droit d'usage. Pour la loi étatique, seul le titre foncier enregistré fait foi — B est juridiquement fondé. Le conflit est donc un conflit de \emph{légitimités}, non seulement de terrains.
\item \textbf{Mécanisme traditionnel} : convocation d'une palabre réunissant les deux chefs, les notables et les sages des deux lignages ; écoute des mémoires généalogiques sur l'histoire du terrain ; recherche d'un compromis (usage partagé, compensation, bail coutumier reconnu).
\item \textbf{Articulation avec la loi étatique} : la commission administrative de conciliation foncière (prévue par la réglementation camerounaise) entérine l'accord issu de la palabre ; le titre de B peut être aménagé par une servitude ou une cession partielle au profit de A. La coutume fournit la \emph{reconnaissance mutuelle}, la loi fournit la \emph{sécurité juridique}.
\end{enumerate}
\textbf{Conclusion philosophique} : le vivre-ensemble n'a pas consisté à faire taire une norme au profit de l'autre, mais à gérer pacifiquement la différence des légitimités — exactement ce que l'ontologie relationnelle prescrit : la personne (ici, chaque village) ne se réalise qu'en maintenant le lien avec l'autre.
\end{exoresolu}

\begin{exercice}[Pour s'entraîner à la dissertation]
Sujet : « Le vivre-ensemble suppose-t-il l'effacement des différences ? » Rédigez l'introduction (accroche, définition des termes, problématique, annonce du plan) en mobilisant l'Ubuntu, le Préambule de la Constitution de 1996 et l'exemple du njangi.
\end{exercice}

\begin{corrige}[Exercice — éléments de correction]
\textbf{Accroche} : le Cameroun, « l'Afrique en miniature », compte plus de 250 groupes ethniques. \textbf{Définitions} : vivre-ensemble (valeur constitutionnelle, Préambule de 1996) ; différence (appartenance culturelle, linguistique). \textbf{Problème} : si la cohésion exige le commun, faut-il dissoudre les particularismes ? \textbf{Problématique} : le vivre-ensemble est-il l'uniformisation ou la gestion pacifique de la différence ? \textbf{Plan annoncé} : I. La tentation de l'uniformité (l'unité semble exiger le même) ; II. Mais l'uniformité détruit le lien (Ubuntu : la personne est relation, or la relation suppose l'altérité ; le njangi unit des membres différents sans les confondre) ; III. Le vivre-ensemble comme unité composée (Charte africaine, art. 28 : tolérance et dialogue).
\end{corrige}

\begin{aretenir}[L'essentiel de la section]
\begin{itemize}
\item L'ontologie africaine est \cle{relationnelle} : « Je suis parce que nous sommes » (Ubuntu), contrepoint du cogito cartésien.
\item La personne est un nœud de relations : famille, cité, ancêtres, cosmos ; le droit est subsidié au devoir (Wiredu).
\item Le vivre-ensemble a des fondements camerounais (njangi, mvett, dina, Ngondo) et un statut constitutionnel (Préambule, 1996).
\item La Charte africaine (1981, art. 27-29) énonce des devoirs corrélatifs des droits.
\item Menaces : tribalisme, haine numérique, crises sécuritaires, corruption. Réponses : ECM, bilinguisme, décentralisation, CNPBM.
\item Formule à retenir : l'unité n'est pas l'uniformité ; le vivre-ensemble est la gestion pacifique de la différence.
\end{itemize}
\end{aretenir}

\coursec{Gouvernance : modèles traditionnels, État moderne et défis démocratiques}

Dans la section précédente, nous avons vu que le \emph{vivre-ensemble} africain repose sur une ontologie relationnelle : « Je suis parce que nous sommes ». Mais une communauté unie par des liens de parenté, de solidarité et de rites ne suffit pas : il faut encore organiser le pouvoir, trancher les conflits, répartir les ressources. Autrement dit, il faut \cle{gouverner}. La question qui structure cette section est donc la suivante : comment le Cameroun est-il passé des pouvoirs traditionnels à l'État moderne, et pourquoi la gouvernance démocratique y demeure-t-elle un défi ?

\courssub{Les modèles traditionnels de gouvernance}

\begin{definition}[Gouvernance]
La \cle{gouvernance} désigne l'ensemble des processus par lesquels une société prend des décisions collectives, exerce le pouvoir, contrôle les dirigeants et règle ses conflits. Elle inclut les institutions formelles (État, tribunaux) et informelles (coutumes, sociétés secrètes, conseils de notables).
\end{definition}

L'intuition de départ est simple : bien avant l'arrivée des colonisateurs, l'Afrique — et le Cameroun en particulier — connaissait des formes organisées de pouvoir politique. Prétendre que la colonisation aurait « apporté » la politique à l'Afrique est une erreur historique que la philosophie africaine (Kagame, Nkrumah, Wiredu) a précisément dénoncée.

\textbf{Typologie des pouvoirs traditionnels au Cameroun.} On distingue classiquement trois grands modèles :

\begin{enumerate}
\item \textbf{Les chefferies centralisées} : chez les \cle{Bamoun} (Foumban), les \cle{Bamiléké} (Grassfields) et dans les \cle{Lamidats} du Nord (Maroua, Garoua, Ngaoundéré), le pouvoir est concentré entre les mains d'un roi ou d'un chef sacré. Sa légitimité est \emph{religieuse} : il est le médiateur entre les vivants, les ancêtres et la terre. Il dispose d'une cour, d'une administration, parfois d'une armée.
\item \textbf{Les sociétés à lignages segmentaires} : chez les \cle{Béti}, \cle{Bassa} ou \cle{Douala}, le pouvoir est \emph{diffus}. Il n'y a pas de roi, mais des chefs de lignage et un \emph{conseil des anciens} qui décide par \cle{consensus} après de longues palabres. Le pouvoir y est collégial et tempéré.
\item \textbf{Les sociétés à classes d'âge} : chez les \cle{Kirdi} et les \cle{Mafa} de l'Extrême-Nord, l'autorité s'exerce par groupes d'âge initiés, chaque classe assumant des fonctions (défense, justice, travaux collectifs) selon un roulement générationnel.
\end{enumerate}

\textbf{Les mécanismes de régulation.} Ces sociétés disposaient de véritables « contre-pouvoirs » :

\begin{itemize}
\item Le \cle{Kwifo} (ou \emph{Ngumba}) dans les Grassfields : société régulatrice qui contrôle le chef lui-même, exécute les décisions de justice et peut destituer un souverain indigne.
\item Le \cle{Mbilan} : conseil des notables qui délibère et modère le pouvoir du chef.
\item La \cle{justice restaurative} : contrairement à la justice purement punitive, elle vise la \emph{réparation} du tort et la \emph{réconciliation} des parties (compensation à la famille lésée, rites de purification), car l'objectif est de restaurer l'harmonie du groupe, non d'éliminer le coupable.
\item L'\cle{inviolabilité de la parole donnée} : la parole (\emph{la palabre}) a valeur de contrat sacré ; la trahir, c'est rompre le lien social lui-même.
\end{itemize}

\begin{savaistu}[Une administration africaine précoloniale sophistiquée]
Le Sultanat \cle{Bamoun} de Foumban possède une écriture originale, le \emph{Shümom}, inventée par le Roi \textbf{Njoya} vers 1896. Grâce à elle, le royaume tenait des archives, une cartographie et une pharmacopée écrites : preuve irréfutable d'une administration étatique africaine développée \emph{avant} la colonisation.
\end{savaistu}

\courssub{La rencontre coloniale : dénaturation des légitimités}

La colonisation ne détruit pas seulement les pouvoirs traditionnels : elle les \emph{dénature} en les instrumentalisant.

\begin{itemize}
\item L'\cle{Indirect Rule} britannique (Ouest du Cameroun anglophone) gouverne \emph{à travers} les chefs existants, en les subordonnant à l'autorité coloniale.
\item La politique française d'\cle{Assimilation} puis d'\cle{Association} crée des \emph{chefs de canton} nommés, de simples « auxiliaires de l'administration » chargés de lever l'impôt et de recruter la main-d'œuvre.
\end{itemize}

\textbf{Conséquence philosophique majeure} : le chef n'est plus légitime par le sacré et le consensus de sa communauté, mais par la \emph{nomination administrative}. La légitimité est retournée : le chef ne doit plus compte à son peuple, mais au gouverneur. C'est la naissance d'un pouvoir déconnecté de la reddition des comptes — un héritage lourd pour l'État postcolonial.

\courssub{L'État postcolonial : centralisation et néopatrimonialisme}

À l'indépendance, le Cameroun hérite d'un État centralisé de type \cle{jacobin}. La loi n° 72/4 du 26 août 1972 parachève cette centralisation : les chefs traditionnels deviennent des « auxiliaires de l'administration » placés sous l'autorité des sous-préfets. Le \cle{parti unique} (UNC puis RDPC) fusionne parti et État, et le pouvoir se \emph{personnalise} autour du chef de l'État.

\begin{definition}[Néopatrimonialisme (J.-F. Bayart)]
Système politique dans lequel les dirigeants traitent l'appareil d'État comme leur \emph{domaine privé} : les ressources publiques (postes, marchés, licences) servent à fidéliser des \emph{clientèles personnelles} plutôt qu'à poursuivre l'intérêt général. L'État devient un « gâteau » à partager entre fidèles — ce que Bayart appelle la « politique du ventre ».
\end{definition}

\begin{propriete}[Deux sources de légitimité]
La \cle{légitimité traditionnelle} repose sur le \emph{sacré} (ancêtres, terre, rites) et le \emph{consensus} communautaire ; la \cle{légitimité moderne} repose sur la \emph{légalité} (élections, Constitution) et le \emph{contrat social}. Une gouvernance légitime au Cameroun doit articuler les deux, sans confondre ni opposer artificiellement ces registres.
\end{propriete}

\courssub{La transition démocratique et ses limites}

Sous la pression de la rue et de l'opinion internationale, les années 1990 ouvrent une transition : la \cle{Conférence tripartite} (1991), le retour du \cle{multipartisme} (lois de 1990), puis la \cle{décentralisation} consacrée par la Constitution de 1996 et la loi n° 2004/017 du 22 juillet 2004, qui transfère des compétences aux régions et aux communes.

Pourtant, des blocages persistent :

\begin{itemize}
\item le \cle{clientélisme} : les votes et les nominations s'échangent contre des faveurs plutôt que sur des programmes ;
\item la faible \cle{reddition des comptes} (\emph{accountability}) : les élus et gestionnaires rendent rarement compte aux citoyens ;
\item le \cle{déficit de légitimité} : abstention élevée, défiance envers les institutions.
\end{itemize}

\begin{exemplebox}[La corruption en chiffres]
Selon \emph{Transparency International} (2023), le Cameroun se classe \textbf{140\textsuperscript{e} sur 180} pays, avec un score de \textbf{26/100} à l'indice de perception de la corruption. La CONAC estime le coût de la corruption à plus de \textbf{1 000 milliards de FCFA par an} — soit l'équivalent de plusieurs budgets de ministères entiers, autant de ressources soustraites aux écoles, aux hôpitaux et aux routes.
\end{exemplebox}

\courssub{Gouvernance locale et participation citoyenne}

Face à ces limites, de nouvelles pratiques émergent, qui renouent d'ailleurs avec l'esprit du \emph{Mbilan} et de la palabre :

\begin{itemize}
\item le \cle{budget participatif} : dans certaines communes pilotes, les habitants débattent et priorisent directement une partie des dépenses publiques ;
\item les \cle{Comités de développement villageois} (CVD) : instances locales où la population co-décide des projets (forages, écoles, ponts) ;
\item le rôle des \cle{ONG} et des \cle{diasporas}, qui financent et contrôlent des projets de développement ;
\item la \cle{gouvernance électronique} (e-gouv) : services en ligne (impôts, casier judiciaire) qui réduisent les contacts directs — donc les occasions de corruption.
\end{itemize}

\begin{popfigure}
\begin{center}
\small
\begin{tabularx}{\linewidth}{|l|X|X|X|}
\hline
\rowcolor{popblueL} \textbf{Critère} & \textbf{Modèle traditionnel} & \textbf{État jacobin postcolonial} & \textbf{Gouvernance démocratique décentralisée} \\
\hline
Source de légitimité & Sacré (ancêtres, terre), consensus & Nomination, parti unique, force & Légalité : élections, Constitution \\
\hline
Mode de désignation & Hérédité lignagère, investiture rituelle & Décret, cooptation administrative & Suffrage universel, élections locales \\
\hline
Reddition des comptes & Conseil des anciens, Kwifo (destitution possible) & Très faible : comptes rendus « vers le haut » & Contrôle citoyen, justice, presse, CONAC \\
\hline
Relation gouvernés / gouvernants & Proximité, palabre, réciprocité & Distance, subordination, clientélisme & Contrat social, participation, e-gouv \\
\hline
Gestion des conflits & Justice restaurative, réconciliation & Tribunaux étatiques, répression & Justice de l'État de droit + médiation \\
\hline
\end{tabularx}
\end{center}
\legende{Tableau comparatif des trois modèles de gouvernance au Cameroun. Lecture : aucun modèle n'existe à l'état pur ; la réalité camerounaise actuelle est hybride.}
\end{popfigure}

\begin{attention}[Piège de la dissertation]
Ne pas \emph{essentialiser} la « tradition » comme un bloc statique, naturellement démocratique et harmonieux — les chefferies connaissaient aussi l'autoritarisme et l'esclavage. Ne pas idéaliser non plus l'État moderne comme purement rationnel — il peut être néopatrimonial. Au Cameroun, tradition et modernité sont des \emph{constructions historiques hybrides} : le chef est à la fois gardien du sacré et auxiliaire de l'administration.
\end{attention}

\begin{methode}[Analyser un discours politique]
Pour soumettre un discours politique à l'épreuve critique :
\begin{enumerate}
\item \textbf{Identifier les références} : le locuteur s'appuie-t-il sur la tradition (ancêtres, « notre père », coutume) ou sur la modernité (Constitution, droits, développement) ?
\item \textbf{Repérer le vocabulaire} : parle-t-il du \emph{peuple-enfant} et du \emph{chef-père} (registre patrimonial) ou du \emph{citoyen} et de ses \emph{droits} (registre républicain) ?
\item \textbf{Évaluer la promesse} : promet-il une \emph{redistribution} de faveurs à des fidèles (clientélisme) ou un \emph{développement} au service de l'intérêt général ?
\item \textbf{Conclure} : le discours renforce-t-il la reddition des comptes ou la personnalisation du pouvoir ?
\end{enumerate}
\end{methode}

\begin{exoresolu}[Application : analyser un discours]
\textbf{Énoncé.} Un notable déclare en campagne : « Votez pour moi, je suis votre fils ; votre village aura son forage et je placerai vos enfants à la fonction publique. » Analysez ce discours avec la grille de la méthode, puis jugez sa conformité à une bonne gouvernance.
\tcblower
\textbf{Solution.}
1. \emph{Références} : registre traditionnel et patrimonial (« votre fils », lien de parenté villageoise).
2. \emph{Vocabulaire} : relation père/fils, don/contre-don, et non citoyen/droits.
3. \emph{Promesse} : redistribution ciblée (un forage pour « son » village, des postes pour « ses » enfants) et non politique publique universelle.
4. \emph{Jugement} : ce discours relève du \cle{clientélisme} et du \cle{néopatrimonialisme} (Bayart) : les ressources publiques (forage, emplois d'État) sont traitées comme des biens personnels à distribuer contre des voix. Il viole le principe d'égal accès aux services publics et affaiblit la reddition des comptes. Une bonne gouvernance exigerait que le forage réponde à un plan communal voté (budget participatif) et que les recrutements soient ouverts à tous par concours.
\end{exoresolu}

\begin{aretenir}[L'essentiel de la section]
\begin{itemize}
\item Le Cameroun précolonial connaissait des pouvoirs organisés : chefferies sacrées centralisées (Bamoun, Bamiléké, Lamidats), sociétés segmentaires à conseil des anciens (Béti, Bassa, Douala), classes d'âge (Kirdi, Mafa), régulées par le Kwifo, le Mbilan et la justice restaurative.
\item La colonisation a dénaturé ces légitimités en transformant les chefs en auxiliaires de l'administration.
\item L'État postcolonial, centralisé et personnalisé, a dérivé vers le \cle{néopatrimonialisme} (Bayart).
\item La transition démocratique (multipartisme, décentralisation de 2004) reste freinée par le clientélisme et le déficit de reddition des comptes.
\item L'avenir passe par la participation citoyenne : budgets participatifs, CVD, ONG, diasporas et e-gouvernance.
\end{itemize}
\end{aretenir}

\begin{exercice}[Pour s'entraîner à la dissertation]
Sujet : \emph{« La gouvernance démocratique au Cameroun doit-elle rompre avec les pouvoirs traditionnels ou s'appuyer sur eux ? »} Rédigez une introduction problématisée (5 à 8 lignes) en mobilisant la distinction entre légitimité traditionnelle et légitimité moderne, puis un plan en deux parties.
\end{exercice}

\begin{corrige}[Exercice 1]
\emph{Introduction attendue} : amorce (toute société doit organiser le pouvoir) ; présentation du problème (le Cameroun hérite de deux légitimités : sacrée/consensuelle et légale/électorale) ; problématique (faut-il opposer ces légitimités ou les articuler ?) ; annonce du plan. \emph{Plan possible} : I. Les pouvoirs traditionnels semblent incompatibles avec la démocratie (hérédité, sacré, absence d'élection) ; II. Pourtant, ils offrent des ressources pour une gouvernance participative (consensus, palabre, justice restaurative, Kwifo comme contre-pouvoir) qu'une démocratie décentralisée peut intégrer. Conclusion ouverte : la bonne gouvernance camerounaise sera hybride, articulant légalité moderne et sagesse traditionnelle.
\end{corrige}

\coursec{Synthèse appliquée : Articuler philosophie, vivre-ensemble et gouvernance au Cameroun}

Après avoir étudié les modèles traditionnels de gouvernance, l'État moderne et les défis démocratiques (section 3), il nous reste à accomplir le geste philosophique par excellence : la \cle{synthèse}. Comment la philosophie africaine, née d'une histoire de résistance et d'affirmation, peut-elle devenir un outil pratique pour penser le \cle{vivre-ensemble} et transformer la \cle{gouvernance} au Cameroun ? Cette dernière section montre que la philosophie n'est pas un savoir abstrait : elle est une force d'action au service de la cité.

\courssub{L'Unité dans la diversité : la devise nationale comme synthèse philosophique}

Le Cameroun compte plus de 250 groupes ethniques, deux grandes zones linguistiques (francophone et anglophone) et plusieurs religions. Comment faire \emph{un} avec tant de diversité ? La devise nationale \og Paix -- Travail -- Patrie \fg{} peut être lue comme une véritable synthèse philosophique en trois termes :

\begin{itemize}
\item \textbf{Paix} : l'\cle{unité ontologique}. Avant d'être un programme politique, la paix affirme que tous les Camerounais partagent la même humanité et la même dignité. Elle renvoie à l'\emph{Ubuntu} : \og je suis parce que nous sommes \fg.
\item \textbf{Travail} : l'\cle{effort collectif}. L'unité ne se décrète pas, elle se construit par la coopération, comme dans les traditions de travail communautaire (les \emph{tontines}, le travail champêtre collectif).
\item \textbf{Patrie} : l'\cle{appartenance commune}. Par-delà les appartenances particulières (ethnie, région, religion), il existe un bien commun supérieur : la nation.
\end{itemize}

\begin{aretenir}[La devise comme philosophie politique]
\og Paix -- Travail -- Patrie \fg{} articule les trois dimensions de toute vie commune : l'être ensemble (paix), le faire ensemble (travail) et le destin partagé (patrie). Une devise n'est pas un slogan : elle est le résumé d'un contrat social.
\end{aretenir}

\courssub{La décentralisation : traduction institutionnelle du vivre-ensemble}

La Constitution de 1996 fait du Cameroun un État \cle{décentralisé}, organisé en régions et en communes. Ce choix institutionnel possède un fondement philosophique précis : le principe de \cle{subsidiarité}.

\begin{definition}[Subsidiarité]
Principe selon lequel l'autorité centrale ne doit intervenir que si l'échelon local ne peut pas assumer lui-même une compétence. C'est le fondement éthique de la décentralisation : ce qui peut être décidé au plus près des citoyens doit l'être.
\end{definition}

La subsidiarité traduit concrètement le vivre-ensemble : elle fait confiance aux communautés locales pour gérer leurs \cle{biens communs} (l'eau, la forêt, le marché). Au Cameroun, les Comités Villageois de Développement (CVD) illustrent cette logique : les habitants décident ensemble des priorités (forage, salle de classe, piste rurale) et participent à leur réalisation. C'est la sagesse des palabres traditionnelles transposée dans l'administration moderne.

\begin{exemplebox}[L'abstention, symptôme d'une crise du vivre-ensemble]
Lors du double scrutin législatif et municipal de février 2020, le taux de participation annoncé par ELECAM était d'environ 45 \%. Plus d'un électeur inscrit sur deux n'a pas voté pour choisir ses édiles locaux, alors même que la commune est l'échelon le plus proche du quotidien. Ce chiffre révèle un défi majeur : réinvestir la \emph{chose publique} (\emph{res publica}). L'éducation philosophique à la citoyenneté est une réponse : comprendre que la politique n'est pas l'affaire des autres, mais notre affaire commune.
\end{exemplebox}

\courssub{Justice transitionnelle et réconciliation : l'apport spécifique de l'Afrique}

Les crises que traverse le Cameroun (crise anglophone dans le Nord-Ouest et le Sud-Ouest, insécurité dans l'Extrême-Nord face à Boko Haram) posent une question philosophique redoutable : comment reconstruire le vivre-ensemble après la violence ?

Deux modèles de justice s'affrontent :

\begin{center}
\begin{tabularx}{\linewidth}{|l|X|X|}
\hline
\rowcolor{popblueL} \textbf{Critère} & \textbf{Justice pénale (internationale)} & \textbf{Justice africaine de réconciliation} \\
\hline
Finalité & Punir le coupable & Réparer le lien social brisé \\
\hline
Acteur central & Le juge, l'État & La communauté, les victimes et les coupables \\
\hline
Moyen & La sanction, la peine & Le pardon, la réparation, la réintégration \\
\hline
Horizon & La rétribution du passé & L'avenir commun \\
\hline
\end{tabularx}
\end{center}

La philosophie africaine apporte ici une contribution originale. Les traditions de palabre, les rites de réconciliation et la philosophie de l'\emph{Ubuntu} montrent qu'une société ne guérit pas seulement en punissant, mais en réintégrant l'offensé et l'offenseur dans la communauté. L'exemple sud-africain de la Commission Vérité et Réconciliation (présidée par Desmond Tutu après l'apartheid) en est l'illustration célèbre : la vérité publique et le pardon ont été préférés, autant que possible, aux procès, pour éviter la guerre civile.

\begin{attention}[Ne pas opposer les deux modèles de façon simpliste]
La justice de réconciliation ne dispense pas de toute sanction : les crimes graves (massacres, crimes de guerre) relèvent de la justice pénale. Le pardon ne signifie pas l'impunité. La synthèse consiste à articuler vérité, justice et réparation : punir les crimes imprescriptibles, mais réintégrer ceux qui peuvent l'être.
\end{attention}

\courssub{L'éthique de la responsabilité : penser les générations futures}

Le philosophe Hans Jonas a formulé le principe d'une \cle{éthique de la responsabilité} : \og Agis de façon que les effets de ton action soient compatibles avec la permanence d'une vie authentiquement humaine sur terre \fg. Autrement dit, notre puissance technique moderne nous rend responsables d'êtres qui ne sont pas encore nés.

Cette éthique rejoint des intuitions profondes de la pensée africaine. Fabien Eboussi Boulaga, philosophe camerounais, a montré que la tradition africaine conçoit la communauté comme incluant les morts, les vivants \emph{et} les enfants à naître : la terre n'est pas une propriété, mais un dépôt confié par les ancêtres pour les descendants. Trois applications concrètes au Cameroun :

\begin{enumerate}
\item \textbf{L'environnement} : la déforestation du bassin du Congo ou l'épuisement des sols hypothèquent l'avenir des générations futures. La cosmovision écologique africaine (la nature comme parente, non comme marchandise) offre des ressources pour penser le changement climatique.
\item \textbf{La dette publique} : emprunter aujourd'hui, c'est faire payer demain. L'éthique de la responsabilité exige que l'endettement serve des investissements utiles (écoles, hôpitaux, routes) et non des dépenses somptuaires.
\item \textbf{L'éducation} : négliger l'école, c'est voler l'avenir des jeunes générations. Investir dans l'éducation est l'acte de responsabilité par excellence.
\end{enumerate}

\courssub{Le philosophe citoyen : produire du savoir pour l'action}

Paulin Hountondji (Bénin) a critiqué la philosophie africaine lorsqu'elle se contente de décrire des traditions : il appelle à une philosophie \emph{critique} et \emph{engagée}, qui \og produit du savoir pour l'action \fg. Le philosophe citoyen n'est pas enfermé dans sa tour d'ivoire : il intervient dans la cité.

Concrètement, cet engagement prend plusieurs formes au Cameroun :

\begin{itemize}
\item les \emph{think-tanks} (laboratoires d'idées) qui éclairent les politiques publiques ;
\item les médias, où l'intellectuel analyse l'actualité et déconstruit les rumeurs et les discours de haine ;
\item la société civile (associations, ONG) qui défend les droits et forme les citoyens ;
\item les partis politiques, où la philosophie peut nourrir les programmes.
\end{itemize}

\begin{attention}[Le piège du \og retour aux sources \fg{} nostalgique]
La philosophie africaine n'est pas une archéologie qui rêverait de revenir à un passé idéal. Chez Kwasi Wiredu ou Achille Mbembe, elle est \emph{critique} (elle examine aussi les défauts des traditions : patriarcat, gérontocratie) et \emph{prospective} (elle prépare l'avenir). Valoriser la tradition, c'est la trier, non la copier.
\end{attention}

\begin{savaistu}[Le consensus contre la tyrannie de la majorité]
Le philosophe ghanéen Kwasi Wiredu propose le \cle{consensus} comme alternative au vote majoritaire (\emph{majority rule}). Dans les palabres akan, on délibère jusqu'à ce que tous acceptent la décision, afin d'éviter qu'une majorité, même d'une seule voix, écrase durablement une minorité. Une piste précieuse pour des sociétés multiethniques comme le Cameroun, où le vote majoritaire peut cristalliser les clivages.
\end{savaistu}

\courssub{Perspectives : philosophie africaine face aux défis du XXI\textsuperscript{e} siècle}

La philosophie africaine ne se contente pas de penser le passé : elle affronte les défis nouveaux.

\textbf{Intelligence artificielle.} Les algorithmes qui filtrent l'information, recrutent ou surveillent posent des questions éthiques : qui possède nos données ? Les biais algorithmiques (programmes entraînés sur des données occidentales qui méconnaissent les réalités africaines) exigent une éthique des données africaine, soucieuse de dignité et de justice.

\textbf{Changement climatique.} L'Afrique, qui pollue le moins, subit le plus les effets du réchauffement (sécheresses dans l'Extrême-Nord, avancée du désert). La cosmovision écologique africaine -- la terre comme mère commune et non comme ressource à exploiter -- peut inspirer une écologie responsable, en dialogue avec l'éthique de Jonas.

\courssub{Carte mentale de synthèse}

\begin{popfigure}
\begin{center}
\begin{tikzpicture}[
  mindmap,
  grow cyclic,
  every node/.style={concept, text=white},
  root concept/.append style={concept color=popink, font=\large\bfseries, minimum size=3.2cm},
  level 1/.append style={level distance=4.6cm, sibling angle=120, font=\small\bfseries},
  level 2/.append style={level distance=2.8cm, sibling angle=45, font=\scriptsize},
  concept connection/.append style={color=popmuted}
]
\node[root concept]{Philosophie africaine au Cameroun}
  child[concept color=popblue]{ node {Histoire et courants}
    child[concept color=popblue!70]{ node {Ethnophilosophie (Tempels)} }
    child[concept color=popblue!70]{ node {Critique (Hountondji, Eboussi Boulaga)} }
    child[concept color=popblue!70]{ node {Sagesse (Oruka)} }
  }
  child[concept color=popgreen]{ node {Vivre-ensemble}
    child[concept color=popgreen!70]{ node {Ubuntu, consensus (Wiredu)} }
    child[concept color=popgreen!70]{ node {Palabre, réconciliation} }
    child[concept color=popgreen!70]{ node {Paix -- Travail -- Patrie} }
  }
  child[concept color=poporange]{ node {Gouvernance}
    child[concept color=poporange!70]{ node {Décentralisation, subsidiarité} }
    child[concept color=poporange!70]{ node {État de droit, démocratie} }
    child[concept color=poporange!70]{ node {Défis : abstention, crises} }
  };
\end{tikzpicture}
\end{center}
\legende{Carte mentale de synthèse. Trois branches principales (histoire et courants ; vivre-ensemble ; gouvernance) sont traversées par trois liens transversaux : l'\emph{éthique de la responsabilité} (Jonas, Eboussi Boulaga), la \emph{décentralisation} (subsidiarité, CVD) et l'\emph{éducation} philosophique à la citoyenneté.}
\end{popfigure}

\courssub{Méthode et entraînement à la dissertation}

\begin{methode}[Dissertation type Bac : \og La philosophie africaine peut-elle fonder la démocratie ? \fg]
\begin{enumerate}
\item \textbf{Expliciter les concepts} : définir la philosophie africaine (pensée critique enracinée dans les cultures africaines) et la démocratie (pouvoir du peuple, État de droit, droits individuels). Dégager la tension : des valeurs communautaires peuvent-elles fonder un régime fondé sur l'individu citoyen ?
\item \textbf{Thèse -- les apports} : le consensus (Wiredu) humanise la décision ; la palabre institue la délibération publique ; l'éthique de la responsabilité donne une légitimité endogène au pouvoir ; l'Ubuntu fonde la dignité de tous.
\item \textbf{Antithèse -- les limites} : certaines pratiques traditionnelles contredisent la démocratie moderne : patriarcat (exclusion des femmes), gérontocratie (exclusion des jeunes), absence de droits individuels formels, risque que le consensus étouffe le dissensus.
\item \textbf{Synthèse -- l'hybridation nécessaire} : ni importation pure du modèle occidental, ni retour nostalgique aux traditions, mais une démocratie \emph{hybride} : État de droit et droits individuels, enrichis des valeurs endogènes (consensus, responsabilité, communauté).
\end{enumerate}
\end{methode}

\begin{exoresolu}[Commentaire d'une situation concrète]
\textbf{Énoncé.} Dans un village de l'Ouest, le chef convoque une palabre pour décider de l'emplacement d'un forage. Après trois heures de discussion, tous les notables acceptent un site, mais les femmes et les jeunes, qui puisent l'eau quotidiennement, n'ont pas été consultés. Analysez cette situation à la lumière de la philosophie africaine et de la démocratie.
\tcblower
\textbf{Solution détaillée.} 1) \emph{Ce qui est conforme à la sagesse africaine} : la palabre réalise l'idéal de délibération publique et de consensus valorisé par Wiredu ; la décision n'est pas imposée d'en haut (subsidiarité respectée). 2) \emph{Ce qui fait problème} : le consensus est incomplet, car il exclut les premières concernées (femmes, jeunes). Il s'agit du défaut classique du patriarcat et de la gérontocratie, critiqués par les philosophes africains eux-mêmes. 3) \emph{Conclusion} : un authentique consensus exige l'inclusion de tous les intéressés. La philosophie africaine, loin de justifier l'exclusion, fournit elle-même le critère de sa critique : l'Ubuntu affirme la dignité égale de chacun. La démocratie locale doit donc hybrider palabre et droits égaux.
\end{exoresolu}

\begin{exercice}[Pour s'entraîner]
1. Rédigez l'introduction d'une dissertation sur le sujet : \og Faut-il punir ou pardonner pour réconcilier une société ? \fg (accroche, définitions, problématique, annonce du plan).
2. Expliquez en dix lignes comment le principe de subsidiarité peut améliorer la gestion d'un bien commun (eau, forêt ou marché) dans votre localité.
3. À l'aide d'un exemple camerounais, montrez que l'éthique de la responsabilité (Jonas) concerne les décisions publiques actuelles.
\end{exercice}

\begin{corrige}[Exercice 2]
Le principe de subsidiarité veut que la compétence soit exercée par l'échelon le plus proche des citoyens capable de l'assumer. Pour un point d'eau villageois, le comité local (CVD) connaît les besoins réels, les lieux de source et les usages ; il peut donc décider de l'emplacement, organiser l'entretien et fixer de petites contributions. L'État central n'intervient que pour ce qui dépasse les capacités locales : forage profond, normes sanitaires, financement lourd. Ainsi la décision est plus juste (prise par les concernés), plus efficace (adaptée au terrain) et plus responsabilisante : en gérant eux-mêmes leur bien commun, les habitants deviennent des citoyens actifs, ce qui illustre concrètement le vivre-ensemble.
\end{corrige}

\newpage

\coursec{Activité d'intégration}

\coursec{La Philosophie africaine}

\courssub{Genèse et historicité de la philosophie africaine}

\begin{definition}[Philosophie africaine]
Ensemble des réflexions critiques, rationnelles et argumentées produites par des Africains ou sur l'Afrique, portant sur les problèmes fondamentaux de l'existence (être, connaissance, morale, politique, esthétique), en dialogue avec la tradition orale et l'histoire coloniale.
\end{definition}

\textbf{Le débat sur l'existence (années 1940-1970).} La question « Y a-t-il une philosophie africaine ? » structure le champ. Trois positions s'affrontent :
\begin{itemize}
\item \textbf{L'ethnophilosophie} (Tempels \emph{La Philosophie bantoue}, Kagamé \emph{La Philosophie bantu-rwandaise}, Mbiti) : la philosophie africaine est implicite dans les langues, mythes, institutions ; elle est communautaire, anonyme, unanime. Risque : figer la pensée dans un « unanimisme » ahistorique.
\item \textbf{Le nationalisme idéologique} (Senghor, Nkrumah, Nyerere, Touré) : la philosophie sert la décolonisation et la construction nationale (Négritude, Ujamaa, Consciencisme). Risque : instrumentalisation politique.
\item \textbf{L'école universaliste / critique} (Hountondji \emph{Sur la philosophie africaine}, Wiredu, Bodunrin, Appiah, Mudimbe) : la philosophie est une activité individuelle, critique, argumentée, soumise aux normes de rationalité universelle. Il n'y a pas « une » philosophie africaine, mais des philosophes africains qui font de la philosophie.
\end{itemize}

\begin{aretenir}[Repères historiographiques]
\begin{itemize}
\item \textbf{Années 1970-1990 :} Professionnalisation (départements universitaires, revues, associations). Hountondji dénonce l'ethnophilosophie comme « mythe » ; Wiredu propose le « décolonisation conceptuelle » (libérer la pensée des catégories coloniales).
\item \textbf{Depuis les années 1990 :} Tournant politique et éthique (Mbembe, Mudimbe, Oyewumi). Questions : subjectivité, genre, postcolonialité, mondialisation, \cle{vivre-ensemble}, \cle{gouvernance}.
\item \textbf{Au Cameroun :} Figures majeures : Fabien Eboussi Boulaga (critique de l'aliénation, éthique de la libération), Patrice Nganang, Jean-Godefroy Bidima (philosophie de la parole, palabre), Pierre Ela (théologie de la libération, éthique sociale).
\end{itemize}
\end{aretenir}

\begin{attention}[Piège épistémologique]
Ne pas confondre \emph{philosophie africaine} (activité critique présente) et \emph{philosophie de l'Afrique} (discours sur l'Afrique, souvent produit par des non-Africains). Ne pas réduire la philosophie africaine à la sagesse des proverbes : le proverbe est matière à philosopher, non philosophie elle-même (Hountondji).
\end{attention}

\courssub{Le « Vivre-ensemble » : fondements ontologiques et enjeux citoyens}

\begin{definition}[Vivre-ensemble]
Art politique et éthique de partager un espace commun par le dialogue, la reconnaissance mutuelle, la gestion pacifique des conflits et la construction d'un destin partagé, sans nier les différences.
\end{definition}

\textbf{Fondements ontologiques.}
\begin{itemize}
\item \textbf{Ubuntu (Afrique australe) / Communautarisme modéré (Gyekye, Mbiti, Menkiti) :} « \emph{Umuntu ngumuntu ngabantu} » (Je suis parce que nous sommes). La personne se constitue dans la relation. \cle{Wiredu} précise : le communautarisme ne nie pas l'individu, mais le définit par ses obligations relationnelles.
\item \textbf{La personne comme charge et projet (Eboussi Boulaga) :} L'humain n'est pas donné, il se fait par l'assumption de responsabilités envers autrui et la cité.
\item \textbf{La palabre (Bidima, Fall) :} Procédure de parole partagée, temps long, écoute, recherche du consensus. Ce n'est pas une simple coutume, mais une \emph{technique de démocratie} (Fall).
\end{itemize}

\textbf{Enjeux citoyens au Cameroun.}
\begin{itemize}
\item \textbf{Du sujet au citoyen :} Passer de l'allégeance au chef/ethnie à l'allégeance à la Constitution et aux lois (patriotisme constitutionnel).
\item \textbf{Pluralisme juridique :} Articuler droit moderne (Code civil, Ord. 74-1) et droit coutumier (lois non écrites, autorités traditionnelles). La loi n° 2011/011 sur la décentralisation renforce le rôle des chefferies comme auxiliaires de l'administration.
\item \textbf{Défis :} Tribalisme, népotisme, corruption, exclusion des jeunes et des femmes, crises anglophones, insécurité (Boko Haram, NOU). Le vivre-ensemble exige \cle{réparation} (justice transitionnelle) et non seulement répression.
\end{itemize}

\begin{exemplebox}[Exemple camerounais : La « Commission Nationale pour la Promotion du Bilinguisme et du Multiculturalisme » (CNPBM)]
Créée en 2017, elle illustre une tentative institutionnelle de vivre-ensemble : traiter les griefs, promouvoir l'équité linguistique/culturelle, prévenir les conflits. Son efficacité dépend de l'indépendance réelle et de la suite donnée à ses recommandations.
\end{exemplebox}

\courssub{Gouvernance : modèles traditionnels, État moderne et défis démocratiques}

\begin{definition}[Gouvernance]
Manière d'exercer le pouvoir (politique, économique, administratif) pour gérer les affaires communes. Elle se juge à l'aune de principes : \cle{légitimité}, \cle{redevabilité} (accountability), \cle{transparence}, \cle{participation}, \cle{efficacité}, \cle{équité}, \cle{État de droit}.
\end{definition}

\textbf{Modèles traditionnels (chefferies, palais, sultanats).}
\begin{itemize}
\item \textbf{Structure :} Chef (Fo, Sultan, Lamido, Chef de village) + Conseil des Anciens/Notables + Sociétés secrètes/initiatiques (régulation).
\item \textbf{Principes :} Le chef est \emph{premier serviteur} et \emph{gardien} (terres, ancêtres, coutumes), non propriétaire absolu. Décision par \cle{consensus} (palabre) et non vote majoritaire (Wiredu : le consensus protège la minorité).
\item \textbf{Contrôles coutumiers :} Droit de remontrance, destitution rituelle, interdits (totems, serments). \emph{Exemple :} Chez les Bamiléké, le \emph{Kamveu} (conseil des neuf notables) peut destituer le Fo.
\end{itemize}

\textbf{État moderne et cadre légal camerounais.}
\begin{itemize}
\item \textbf{Constitution 1996 (rév. 2008) :} Séparation des pouvoirs, multipartisme, décentralisation (Collectivités Territoriales Décentralisées - CTD), Cour Suprême, Conseil Constitutionnel.
\item \textbf{Chefferies traditionnelles :} Reconnues comme « auxiliaires de l'administration » (Décret 77/245, Loi 2011/011). Elles perçoivent des taxes, gèrent les conflits fonciers coutumiers, mais leur pouvoir est encadré par l'administration (Sous-Préfet, Préfet, Gouverneur).
\item \textbf{Foncier :} Ordonnance 74-1 (domaine national, immatriculation = seul titre de propriété) vs Ordonnance 74-2 (domanialité de l'État) vs Coutume (terre communautaire, inaliénable). Source majeure de contentieux.
\end{itemize}

\begin{propriete}[Critères de bonne gouvernance (PNUD / UA / MINESEC)]
\begin{itemize}
\item \textbf{Participation :} Tous (hommes, femmes, jeunes, minorités) ont voix au chapitre.
\item \textbf{État de droit :} Règles claires, publiques, appliquées équitablement ; justice indépendante.
\item \textbf{Transparence :} Accès à l'information, budgets publics lisibles.
\item \textbf{Reddevabilité :} Les gouvernants rendent des comptes aux gouvernés (élections, reddition de comptes, pétitions).
\item \textbf{Équité / Inclusion :} Pas de discrimination ; protection des vulnérables.
\end{itemize}
\end{propriete}

\begin{attention}[Tensions structurelles]
\begin{itemize}
\item \textbf{Légalité vs Légitimité :} Un chef élu par l'administration peut manquer de légitimité coutumière ; un chef coutumier légitime peut violer la loi (droits des femmes, vente de terres).
\item \textbf{Clientélisme vs Reddevabilité :} La « politique du ventre » (Bayart) transforme la chose publique en ressource privée.
\item \textbf{Jeunesse exclue :} Majorité démographique, minorité décisionnelle. Le vivre-ensemble impose leur intégration effective (conseils de jeunesse, budgets participatifs).
\end{itemize}
\end{attention}

\courssub{Synthèse appliquée : Articuler philosophie, vivre-ensemble et gouvernance au Cameroun}

\textbf{Objectifs de l'activité}

À l'issue de cet atelier de simulation intitulé « La Palabre de Mvomeka'a », l'élève de Terminale technique doit être capable de :

\begin{itemize}
\item \cle{Appliquer} les notions de l'unité (philosophie africaine, vivre-ensemble, gouvernance) à une situation concrète de la vie courante : un conflit foncier opposant légitimité coutumière et légitimité étatique ;
\item \cle{Mobiliser} les concepts philosophiques de \emph{consensus} (Wiredu), d'\emph{Ubuntu} (humanité relationnelle), de \cle{réparation} et de critique de l'ethnophilosophie (Hountondji) pour construire une argumentation rigoureuse ;
\item \cle{Rédiger et justifier} une démarche, une réponse ou une conclusion, sous la forme d'un « Arrêté de médiation » argumenté philosophiquement ;
\item \cle{Distinguer} les sources de légitimité (coutume, loi, intérêt économique, mobilisation populaire) et évaluer leur validité respective ;
\item \cle{Adopter} une posture de médiateur : écoute, impartialité, recherche d'une solution équitable plutôt que d'un vainqueur.
\end{itemize}

\textbf{Situation}

\textbf{Le conflit foncier de Mvomeka'a.}

Mvomeka'a est un village d'environ \num{3200} habitants situé dans la région du Sud du Cameroun, à une quarantaine de kilomètres d'Ebolowa. Le village dispose d'un terrain communal de \SI{12}{hectares}, situé en bordure de la route nationale, sur lequel les familles pratiquent depuis des générations la culture du manioc, du plantain et du cacao. Ce terrain n'a jamais fait l'objet d'une immatriculation au nom de la communauté : il est régi par la \cle{coutume}, c'est-à-dire par les règles non écrites transmises oralement, dont le \emph{Conseil des Anciens} est le gardien.

En janvier 2026, Sa Majesté Ngoa Essomba, chef du village, signe seul — sans consulter le Conseil des Anciens ni les familles utilisatrices — une convention de vente du terrain à la société immobilière « Horizon Développement SARL » pour un montant de 18 millions de francs CFA, soit \num{1500000} FCFA par hectare. Le promoteur entend y construire un lotissement de 60 parcelles résidentielles destinées à des cadres d'Ebolowa. Le chef justifie sa décision par la nécessité de financer la réhabilitation de l'école primaire du village (devis estimé : 9 millions de francs CFA) et affirme que « le chef est le propriétaire naturel des terres de sa chefferie ».

Apprenant la nouvelle, un collectif de 45 jeunes du village, dont 12 exploitent le terrain avec leurs familles, organise une manifestation pacifique devant la chefferie. Ils brandissent des pancartes : « Notre terre n'est pas à vendre », « Le chef n'est pas le village ». Trois d'entre eux sont brièvement interpellés par la gendarmerie, puis relâchés. Alertée, la Sous-Préfecture \emph{suspend} la procédure d'immatriculation du titre foncier demandée par le promoteur, le temps qu'une enquête administrative soit menée.

Le village est désormais divisé : les partisans du chef (environ un tiers des habitants, selon les notables) invoquent l'autorité traditionnelle ; les jeunes et les familles concernées invoquent le droit coutumier bafoué ; le promoteur menace d'engager des poursuites pour « rupture abusive de contrat » et réclame le remboursement de l'avance de 6 millions de francs CFA déjà versée.

Face à l'escalade, le Sous-Préfet propose l'organisation d'une \cle{palabre} de médiation, selon la tradition africaine de la parole partagée, pour trouver une issue équitable. C'est cette palabre que la classe va simuler.

\textbf{Documents remis aux participants :}

\begin{itemize}
\item \textbf{Document 1 — Extrait de l'Ordonnance n° 74-1 du 6 juillet 1974} fixant le régime foncier au Cameroun : « Les terres non immatriculées et non occupées de façon effective font partie du domaine national. L'immatriculation est le seul mode de constitution du droit de propriété foncière. »
\item \textbf{Document 2 — Coutume locale} (relevé oral consigné par les notables) : « La terre du village appartient à la communauté entière, vivants, ancêtres et générations futures. Le chef en est le gardien, non le propriétaire. Aucune cession de terre communautaire n'est valable sans l'accord du Conseil des Anciens réuni en palabre et l'assentiment des familles qui l'exploitent. »
\item \textbf{Document 3 — Témoignages} : celui du chef (« J'ai agi pour l'école ; un chef décide pour le bien de tous ») ; celui d'une jeune cultivatrice (« Ce champ nourrit ma famille depuis ma grand-mère ») ; celui du promoteur (« J'ai signé en bonne foi, j'ai investi, je veux mon terrain ou mon argent ») ; celui du Sous-Préfet (« La loi exige l'immatriculation ; mais la paix sociale exige le dialogue »).
\end{itemize}

\textbf{Consignes}

\textbf{Phase 1 — Lecture et analyse des documents (15 minutes).}

\begin{enumerate}
\item Identifiez, pour chaque document, la \cle{source de légitimité} invoquée (loi étatique, coutume, intérêt économique, autorité traditionnelle) et formulez en une phrase la thèse qu'il défend.
\item Relevez la contradiction fondamentale entre le Document 1 et le Document 2. À quelle conception de la propriété chacun renvoie-t-il ?
\end{enumerate}

\textbf{Phase 2 — Préparation des arguments par rôle (20 minutes).} La classe est répartie en groupes : le médiateur, le chef et ses partisans, les jeunes et familles, le promoteur, l'administration, et les observateurs.

\begin{enumerate}
\setcounter{enumi}{2}
\item Chaque groupe rédige au moins \textbf{trois arguments} en s'appuyant explicitement sur les auteurs étudiés : \emph{Wiredu} (le consensus comme fondement de la décision légitime), \emph{Ubuntu} (« je suis parce que nous sommes » : la personne définie par la relation communautaire), \emph{Hountondji} (la critique de la pensée africaine réduite à un unanimisme coutumier : la tradition elle-même doit être discutée rationnellement).
\item Les observateurs préparent une grille : usage correct des concepts, qualité de l'argumentation, respect de la procédure de la palabre, équité de la solution proposée.
\end{enumerate}

\textbf{Phase 3 — Simulation de la palabre (30 minutes).}

\begin{enumerate}
\setcounter{enumi}{4}
\item Le médiateur ouvre la parole, donne la parole à chaque partie à tour de rôle, reformule les positions, et conduit le groupe vers un \cle{consensus} : aucune décision n'est valable si une partie essentielle est exclue de la parole.
\end{enumerate}

\textbf{Phase 4 — Rédaction individuelle (15 minutes).}

\begin{enumerate}
\setcounter{enumi}{5}
\item Rédigez individuellement un « Arrêté de médiation » (20 à 30 lignes) qui : (a) rappelle les faits et les prétentions ; (b) tranche sur la validité de la vente en justifiant philosophiquement ; (c) propose une solution équitable pour chaque partie ; (d) conclut sur ce que ce cas révèle du vivre-ensemble et de la gouvernance au Cameroun.
\end{enumerate}

\textbf{Ressources à mobiliser}

\begin{aretenir}[Notions et repères utiles]
\begin{itemize}
\item \textbf{Consensus (Kwasi Wiredu)} : dans la tradition de décision africaine (par exemple ashanti), la décision légitime n'est pas celle de la majorité qui écrase la minorité, mais celle qui résulte d'une délibération poursuivie jusqu'à l'accord de toutes les parties. Le consensus est une exigence rationnelle, non une soumission au groupe.
\item \textbf{Ubuntu} : « Je suis parce que nous sommes. » La personne humaine se réalise dans la relation ; une décision qui détruit le tissu communautaire (exproprier des familles sans parole) est injuste même si elle est légale.
\item \textbf{Critique de l'ethnophilosophie (Paulin Hountondji)} : il n'existe pas une « pensée africaine » unanime et intangible ; la philosophie africaine est un débat critique entre individus. On peut donc discuter la coutume elle-même, et le chef n'a pas le monopole de la raison du village.
\item \textbf{Vivre-ensemble} : art de partager un espace commun par le dialogue, la reconnaissance mutuelle et la réparation des torts.
\item \textbf{Gouvernance} : manière d'exercer le pouvoir ; elle exige légitimité, redevabilité (rendre compte), participation et transparence. Le chef qui décide seul manque à la redevabilité ; l'État qui ignore la coutume manque à l'ancrage social de sa légalité.
\end{itemize}
\end{aretenir}

\textbf{Démarche et corrigé commenté}

\begin{exoresolu}[Corrigé commenté de l'atelier « La Palabre de Mvomeka'a »]
Énoncé : conduire la médiation du conflit foncier de Mvomeka'a et rédiger un arrêté de médiation philosophiquement justifié.
\tcblower
\textbf{Étape 1 — Analyse des légitimités en présence.} Quatre sources de légitimité s'affrontent. (a) La \emph{légitimité étatique} (Document 1) : seule l'immatriculation crée la propriété ; or le terrain n'est pas immatriculé, donc juridiquement le chef ne pouvait pas vendre un droit de propriété qui n'existe pas encore — la vente est fragile devant la loi elle-même. (b) La \emph{légitimité coutumière} (Document 2) : le chef est gardien et non propriétaire ; la cession sans le Conseil des Anciens viole la règle coutumière. (c) La \emph{légitimité économique} du promoteur : il a contracté de bonne foi et avancé 6 millions de francs CFA. (d) La \emph{légitimité populaire} des jeunes : leur manifestation exprime un tort réel (perte des moyens de subsistance de 12 familles), mais la violence ou le désordre ne créent pas de droit.

\textbf{Étape 2 — Diagnostic philosophique.} Le chef invoque une conception absolutiste de l'autorité traditionnelle. Or, avec \emph{Wiredu}, on montre que dans la tradition de la palabre, le pouvoir du chef est précisément limité par l'exigence de consensus : décider seul, c'est trahir la tradition au nom de la tradition. Avec \emph{Ubuntu}, la vente ignore que la terre est le support de la relation communautaire (vivants, ancêtres, générations futures) : 12 familles perdent leur subsistance, le « nous » est brisé. Avec \emph{Hountondji}, on refuse cependant de sacraliser la coutume : c'est par la discussion rationnelle, et non par l'invocation magique de « la tradition », que le village doit trancher — et les jeunes ont droit à la parole autant que les anciens.

\textbf{Étape 3 — Conduite de la palabre.} Le médiateur applique la procédure du consensus : chaque partie expose, est reformulée, puis interrogée. Les questions décisives posées au chef : « Avez-vous consulté ? Pouvez-vous rendre compte de l'usage des 6 millions ? » — elles mettent en lumière le déficit de \emph{redevabilité}, critère central de la bonne gouvernance. Question posée au promoteur : « Pouviez-vous ignorer qu'un terrain non immatriculé et cultivé ne pouvait être vendu par une seule personne ? » — sa bonne foi est relative.

\textbf{Étape 4 — L'Arrêté de médiation (modèle de solution équitable).} (a) \emph{Annulation de la vente} : elle viole à la fois la coutume (absence du Conseil des Anciens) et l'esprit de la loi (pas de droit de propriété constitué). (b) \emph{Réparation} : le chef rembourse les 6 millions au promoteur, si nécessaire par échéancier contrôlé par le Conseil des Anciens et le Sous-Préfet ; le promoteur renonce aux poursuites, sa négligence étant établie. (c) \emph{Mesure d'avenir} : la communauté engage l'immatriculation collective du terrain au nom du village, avec l'appui de l'administration — ainsi la coutume et la loi s'articulent au lieu de s'ignorer. (d) \emph{Financement de l'école} : un comité villageois (chef, notables, jeunes, femmes) recherche des financements transparents (budget participatif, subvention communale), car la fin (l'école) était louable mais ne justifiait pas les moyens. (e) \emph{Réconciliation} : le chef reconnaît publiquement sa faute de procédure ; les jeunes reconnaissent l'autorité de la chefferie réformée. C'est la logique de la \emph{réparation} plutôt que de la punition, caractéristique du vivre-ensemble africain.

\textbf{Étape 5 — Conclusion philosophique.} Ce cas montre que la gouvernance légitime au Cameroun ne peut être ni le rejet de la modernité étatique ni l'écrasement des légitimités coutumières : elle exige leur articulation critique. Le vivre-ensemble n'est pas l'absence de conflit, mais la capacité d'une communauté à transformer le conflit en délibération. La palabre, loin d'être un folklore, apparaît comme une procédure rationnelle de production du consensus.
\end{exoresolu}

\textbf{Bilan de l'activité}

Cet atelier a d'abord permis de vérifier que les notions de l'unité ne sont pas de simples définitions à réciter : les concepts de \cle{consensus}, d'\cle{Ubuntu}, de \cle{réparation} et de \cle{gouvernance} fonctionnent comme des outils d'analyse et de résolution d'un conflit réel, de type foncier, fréquent au Cameroun. L'élève a appris à distinguer les sources de légitimité (loi, coutume, économie, mobilisation) et à ne les accepter qu'après examen critique, suivant la leçon de Hountondji : ni soumission à la tradition, ni soumission à la loi nue, mais délibération raisonnée.

Ensuite, la simulation a mis en évidence la dimension \emph{procédurale} de la justice : une décision juste n'est pas seulement une décision au contenu correct, c'est une décision prise selon une procédure où toutes les parties ont eu la parole — ce qui prépare directement à la dissertation sur le vivre-ensemble et la gouvernance. Enfin, la rédaction de l'« Arrêté de médiation » a entraîné le savoir-faire central de l'épreuve : \cle{rédiger et justifier} une démarche et une conclusion, compétence immédiatement transférable à la dissertation et au commentaire de texte du Baccalauréat technique.

\begin{attention}[Erreurs fréquentes à éviter]
Ne pas opposer caricaturalement « tradition contre modernité » : la coutume elle-même condamne ici la décision solitaire du chef. Ne pas confondre consensus et unanimité imposée : le consensus de Wiredu est le résultat d'un débat rationnel, pas le silence des faibles. Ne pas oublier de justifier chaque point de l'arrêté par un concept ou un auteur du cours : une solution équitable non justifiée philosophiquement ne vaut pas la moitié des points dans la grille critériée.
\end{attention}

\newpage

\coursec{Méthodes et exercices résolus}

\coursec{La Philosophie africaine}

\courssub{Savoir-faire 1 : Appliquer les notions à une situation concrète de la vie courante}

\begin{methode}[Analyser une situation concrète avec les notions de la leçon]
\begin{enumerate}
\item \textbf{Lire attentivement la situation} proposée (un fait de société, un événement, un texte bref) et repérer les éléments essentiels : qui ? quoi ? où ? quel problème ?
\item \textbf{Dégager le problème philosophique} : formuler ce que la situation pose comme question (justice, autorité, identité, vivre-ensemble, pouvoir...).
\item \textbf{Choisir la notion adaptée} de l'unité (philosophie africaine, ubuntu, palaver, gouvernance, État de droit...) et la \textbf{définir brièvement}.
\item \textbf{Établir le lien} : montrer explicitement en quoi la notion éclaire la situation (ce n'est pas la situation qui illustre d'elle-même, c'est vous qui faites le lien).
\item \textbf{Mobiliser un auteur} si possible (Tempels, Kagame, Senghor, Mbiti, Wiredu...) pour appuyer l'analyse.
\item \textbf{Conclure} par une appréciation nuancée : ce que la notion permet de comprendre, et ses limites éventuelles.
\end{enumerate}
\textbf{Formule à retenir :} Situation $\rightarrow$ Problème $\rightarrow$ Notion définie $\rightarrow$ Application $\rightarrow$ Jugement nuancé.
\end{methode}

\begin{exoresolu}[Application : la réunion de crise du village de Mbanga]
Énoncé : Dans un quartier de Mbanga, un différend oppose deux familles au sujet d'une parcelle de terre. Le chef du quartier convoque les anciens, les notables et les deux parties. Après de longues heures de discussions où chacun s'exprime librement, un accord est trouvé : la parcelle est partagée et les deux familles scellent la réconciliation par un repas commun. Analysez cette situation à la lumière des notions de la philosophie africaine.
\tcblower
\textbf{Solution détaillée.}

\textbf{Étape 1 — Repérage des éléments de la situation.} Un conflit foncier entre deux familles ; l'intervention d'une autorité traditionnelle (le chef) ; une instance de discussion (les anciens et notables réunis) ; la parole libre donnée à chacun ; une issue consensuelle ; un rite de réconciliation (le repas commun).

\textbf{Étape 2 — Problème philosophique.} Comment une communauté peut-elle résoudre ses conflits sans recours à la violence ni au tribunal étatique ? C'est le problème du \cle{vivre-ensemble} et de la régulation sociale.

\textbf{Étape 3 — Notions mobilisées et définies.}
\begin{itemize}
\item La \cle{palaver} (ou palabre) : méthode traditionnelle africaine de résolution des conflits par la parole, la discussion prolongée et la recherche du consensus.
\item L'\cle{ubuntu} : « je suis parce que nous sommes » (John Mbiti : \emph{I am because we are}) — la personne se définit par son appartenance à la communauté, ce qui explique que la solution vise la réconciliation et non la victoire d'une partie.
\item La \cle{gouvernance traditionnelle} : le chef et les notables exercent une autorité légitime fondée sur la sagesse et l'expérience.
\end{itemize}

\textbf{Étape 4 — Application.} La convocation de tous les intéressés correspond à la palaver : la parole circule, chacun est entendu. Le partage de la parcelle (et non l'exclusion d'une famille) traduit la logique du consensus : l'objectif n'est pas de désigner un coupable mais de \textbf{réparer le lien social}. Le repas commun scelle symboliquement la restauration de l'harmonie, conformément à l'ubuntu : l'individu n'existe pleinement que dans une communauté apaisée.

\textbf{Étape 5 — Conclusion nuancée.} Cette situation montre que la philosophie africaine n'est pas une spéculation abstraite mais une \emph{sagesse pratique} incarnée dans des institutions. Toutefois, la méthode traditionnelle a des limites : elle peut être lente et doit s'articuler au droit moderne (titres fonciers, tribunaux) lorsque le consensus échoue.
\end{exoresolu}

\courssub{Savoir-faire 2 : Rédiger et justifier une démarche, une réponse ou une conclusion}

\begin{methode}[Construire une réponse argumentée et justifiée]
\begin{enumerate}
\item \textbf{Annoncer la démarche} : en une phrase, dites ce que vous allez faire (« Nous montrerons d'abord..., puis..., enfin... »).
\item \textbf{Définir les termes clés} du sujet avant de répondre : aucune réponse sans définitions.
\item \textbf{Argumenter} : chaque affirmation doit être soutenue par une raison, un exemple ou une référence à un auteur. Structure : \emph{thèse $\rightarrow$ argument $\rightarrow$ exemple $\rightarrow$ mini-conclusion}.
\item \textbf{Envisager l'objection} : une réponse justifiée anticipe le contre-argument et y répond.
\item \textbf{Justifier la conclusion} : la conclusion doit découler logiquement des étapes précédentes, jamais apparaître comme un jugement gratuit. Reprenez la question posée et répondez-y explicitement.
\end{enumerate}
\textbf{Réflexe :} après chaque paragraphe, demandez-vous : « Qu'est-ce qui justifie ce que je viens d'écrire ? »
\end{methode}

\begin{exoresolu}[Réponse justifiée : la philosophie africaine existe-t-elle ?]
Énoncé : Certains philosophes occidentaux (comme Hegel ou Lévy-Bruhl) ont soutenu que l'Afrique n'avait pas de philosophie, car la pensée africaine serait « prélogique » et purement mythique. Rédigez une réponse argumentée et justifiée à cette thèse.
\tcblower
\textbf{Solution détaillée.}

\textbf{Annonce de la démarche.} Nous définirons d'abord la philosophie, puis nous présenterons la thèse négatrice et ses présupposés, ensuite nous la réfuterons par l'histoire et les tendances de la philosophie africaine, avant de conclure.

\textbf{Définition.} La \cle{philosophie} est une démarche rationnelle, critique et réflexive qui interroge les grands problèmes de l'existence (l'être, la connaissance, la morale, la politique) à l'aide d'arguments.

\textbf{Thèse adverse.} Hegel exclut l'Afrique de l'histoire universelle ; Lévy-Bruhl parle de « mentalité prélogique ». Présupposé : la philosophie serait uniquement écrite, individuelle et de type grec.

\textbf{Réfutation argumentée.}
\begin{itemize}
\item \emph{Argument historique} : l'Afrique a connu une pensée réflexive ancienne (sagesse de l'Égypte pharaonique, philosophie éthiopienne de Zera Yacob au XVII\textsuperscript{e} siècle, qui réfléchit de façon critique sur la religion et la raison).
\item \emph{Argument conceptuel} : Placide Tempels (\emph{La Philosophie bantoue}, 1945) a montré que les langues et coutumes bantoues reposent sur une ontologie cohérente : la « force vitale ». Alexis Kagame a démontré que les langues bantoues contiennent des catégories philosophiques comparables à celles d'Aristote.
\item \emph{Argument sur les tendances} : la philosophie africaine contemporaine possède des courants organisés — l'ethnophilosophie (Tempels, Kagame), la philosophie critique (Hountondji, Wiredu, Towa), la philosophie héroïque (Oruka) et la philosophie de la négritude (Senghor, Césaire). Une pensée qui se discute elle-même est bien une philosophie.
\end{itemize}

\textbf{Objection et réponse.} On objectera que la pensée africaine est orale et collective. Mais l'oralité n'exclut pas la rationalité : Socrate n'a rien écrit, et personne ne lui refuse le titre de philosophe.

\textbf{Conclusion justifiée.} Puisque la philosophie est démarche rationnelle et critique, et puisque l'Afrique possède des systèmes de pensée cohérents, des sages et des écoles critiques, la négation de la philosophie africaine repose sur un préjugé ethnocentrique, non sur des arguments. La philosophie africaine existe et s'écrit aujourd'hui dans l'histoire universelle de la philosophie.
\end{exoresolu}

\courssub{Savoir-faire 3 : Construire une problématique à partir d'une notion}

\begin{methode}[Problématiser une notion (vivre-ensemble, gouvernance)]
\begin{enumerate}
\item \textbf{Définir la notion} et repérer sa tension interne : toute notion philosophique contient un problème (ex. vivre-ensemble : comment des hommes différents peuvent-ils former un « nous » ?).
\item \textbf{Formuler la question principale} : une seule question, claire, qui commande tout le devoir.
\item \textbf{Décliner en questions secondaires} (2 ou 3) qui ouvrent les parties : fondements ? conditions ? limites ?
\item \textbf{Vérifier} que les questions ne sont pas des questions de fait (réponse oui/non immédiate) mais de \emph{droit} (elles appellent une justification).
\end{enumerate}
\end{methode}

\begin{exoresolu}[Problématiser : le vivre-ensemble au Cameroun]
Énoncé : Le Cameroun compte plus de 250 groupes ethniques, deux grandes régions linguistiques (francophone et anglophone) et plusieurs religions. Construisez une problématique complète sur le vivre-ensemble à partir de cette réalité.
\tcblower
\textbf{Solution détaillée.}

\textbf{Définition et tension.} Le \cle{vivre-ensemble} désigne la capacité d'individus et de groupes différents à partager un espace commun, des institutions et un destin. Tension interne : l'unité suppose du commun, mais la société camerounaise est faite de diversités ; comment le « un » peut-il naître du « multiple » sans l'écraser ?

\textbf{Question principale.} Sur quels fondements le vivre-ensemble est-il possible dans une société plurielle comme le Cameroun ?

\textbf{Questions secondaires.}
\begin{enumerate}
\item Quels fondements ontologiques (ubuntu, solidarité communautaire) la tradition africaine offre-t-elle au vivre-ensemble ?
\item Quelles institutions politiques (État de droit, justice, école, bilinguisme) peuvent garantir l'unité dans la diversité ?
\item Quels obstacles (tribalisme, exclusion, injustice) menacent le vivre-ensemble, et comment les surmonter ?
\end{enumerate}

\textbf{Vérification.} Ces questions exigent une justification (fondements, conditions, remèdes) et non un simple constat : elles sont bien philosophiques. Elles préparent un plan en trois parties : fondements $\rightarrow$ conditions institutionnelles $\rightarrow$ obstacles et solutions.
\end{exoresolu}

\courssub{Savoir-faire 4 : Comparer deux modèles (gouvernance traditionnelle et État moderne)}

\begin{methode}[Rédiger une comparaison équilibrée]
\begin{enumerate}
\item \textbf{Définir chacun des deux termes} séparément.
\item \textbf{Comparer par critères} (et non en deux blocs juxtaposés) : légitimité, exercice du pouvoir, place du citoyen, limites. Un tableau ou un plan par critères est idéal.
\item \textbf{Dégager convergences et divergences} : ce qui rapproche et ce qui oppose.
\item \textbf{Conclure par une articulation} : montrer comment les deux modèles peuvent se compléter (décentralisation, chefferies reconnues par l'État camerounais).
\end{enumerate}
\textbf{Erreur à éviter :} décrire le modèle A puis le modèle B sans jamais les mettre en relation.
\end{methode}

\begin{exoresolu}[Comparaison : chefferie traditionnelle et État moderne au Cameroun]
Énoncé : Comparez la gouvernance traditionnelle africaine et la gouvernance de l'État moderne, puis montrez comment elles s'articulent au Cameroun.
\tcblower
\textbf{Solution détaillée.}

\textbf{Définitions.} La \cle{gouvernance traditionnelle} désigne l'exercice du pouvoir par les chefs, notables et conseils d'anciens, selon la coutume. La \cle{gouvernance moderne} désigne l'exercice du pouvoir par les institutions de l'État (Constitution, lois, suffrage universel, séparation des pouvoirs).

\textbf{Comparaison par critères.}
\begin{center}
\begin{tabularx}{\linewidth}{|l|X|X|}\hline
\rowcolor{popblueL} \textbf{Critère} & \textbf{Gouvernance traditionnelle} & \textbf{État moderne} \\ \hline
Légitimité & Coutume, lignage, sagesse des anciens & Constitution, élections, droit \\ \hline
Exercice du pouvoir & Conseil des notables, palaver, consensus & Institutions, lois, administration \\ \hline
Place du citoyen & Membre de la communauté, parole au sein du groupe & Citoyen doté de droits individuels \\ \hline
Limites & Peut exclure (femmes, jeunes, étrangers) ; lenteur & Peut être lointain, bureaucratique, corrompu \\ \hline
\end{tabularx}
\end{center}

\textbf{Convergences et divergences.} Les deux visent l'ordre social et la résolution des conflits (convergence de finalité), mais ils divergent sur la source de la légitimité (coutume contre droit) et sur le statut de l'individu (membre contre citoyen).

\textbf{Articulation au Cameroun.} L'État camerounais reconnaît officiellement les chefferies traditionnelles (chefs de 1\textsuperscript{er}, 2\textsuperscript{e} et 3\textsuperscript{e} degrés), qui servent d'auxiliaires de l'administration : le chef mobilise la population, règle les petits litiges, transmet les décisions publiques. La décentralisation (régions, communes) rapproche aussi le pouvoir des citoyens.

\textbf{Conclusion.} Ni rejet pur de la tradition ni repli sur elle : la bonne gouvernance au Cameroun suppose une \emph{synthèse} où les valeurs traditionnelles (consensus, écoute, responsabilité) irriguent un État de droit moderne, transparent et redevable.
\end{exoresolu}

\courssub{Savoir-faire 5 : Rédiger une conclusion de dissertation ou de commentaire}

\begin{methode}[Rédiger et justifier une conclusion]
\begin{enumerate}
\item \textbf{Répondre explicitement} à la question posée dans le sujet (reprendre ses termes).
\item \textbf{Bilan synthétique} : rappeler en deux ou trois phrases le chemin parcouru (les grandes étapes du devoir), sans répéter les exemples.
\item \textbf{Justifier la position adoptée} : la conclusion n'est pas une opinion, elle doit apparaître comme le résultat nécessaire de l'analyse.
\item \textbf{Ouverture raisonnée} (facultative mais valorisée) : poser une nouvelle question qui prolonge le sujet sans le trahir.
\end{enumerate}
\textbf{Interdits :} introduire une notion nouvelle, contredire le développement, conclure par « chacun pense ce qu'il veut ».
\end{methode}

\begin{exoresolu}[Rédiger une conclusion : « La bonne gouvernance est-elle possible en Afrique ? »]
Énoncé : À l'issue d'un devoir ayant montré (I) les défis de la gouvernance en Afrique (corruption, népotisme, coups d'État), (II) les ressources de la tradition (redevabilité du chef devant la communauté, palaver) et (III) les exigences de l'État de droit (transparence, élections, justice), rédigez la conclusion du devoir.
\tcblower
\textbf{Solution détaillée.}

\textbf{Réponse explicite à la question.} La bonne gouvernance est-elle possible en Afrique ? Notre analyse permet de répondre : oui, elle est possible, mais à condition d'être construite et non simplement importée.

\textbf{Bilan synthétique.} Nous avons d'abord reconnu les obstacles réels : la corruption et la confiscation du pouvoir fragilisent les États africains. Nous avons ensuite montré que la tradition africaine n'est pas un obstacle mais une ressource : le chef traditionnel devait rendre compte devant le conseil des anciens, et la palaver institue l'exigence de transparence et de participation. Enfin, l'État de droit moderne apporte les instruments indispensables : Constitution, séparation des pouvoirs, justice indépendante, élections régulières.

\textbf{Justification de la position.} Cette réponse n'est pas un optimisme gratuit : elle découle du constat que les défis identifiés ne tiennent pas à une prétendue « nature » africaine, mais à des pratiques que des institutions solides et des valeurs assumées peuvent corriger. Là où la redevabilité traditionnelle rencontre la transparence moderne, la gouvernance s'améliore.

\textbf{Ouverture raisonnée.} Reste une question que notre devoir n'a pu trancher : quelle place accorder à l'éducation civique des jeunes générations pour que la bonne gouvernance devienne une culture partagée et non une simple règle imposée ?
\end{exoresolu}

\courssub{Savoir-faire 6 : Mobiliser correctement une référence d'auteur}

\begin{methode}[Citer un auteur sans se tromper]
\begin{enumerate}
\item \textbf{Retenir l'essentiel} de chaque auteur : nom, thèse principale, éventuellement titre de l'œuvre. Pas besoin de citation mot à mot.
\item \textbf{Introduire la référence} dans la phrase : « Comme le montre Tempels... », « Selon Mbiti... ».
\item \textbf{Expliquer la référence} : une référence jetée sans explication ne vaut rien ; dites en quoi elle appuie votre argument.
\item \textbf{Vérifier l'adéquation} : l'auteur cité doit réellement défendre la thèse qu'on lui prête.
\end{enumerate}
\textbf{Références clés de l'unité :} Tempels (force vitale bantoue), Kagame (philosophie comparée des Bantous), Senghor (négritude, « raison-étreinte »), Hountondji et Towa (critique de l'ethnophilosophie), Mbiti (communautarisme : « je suis parce que nous sommes »), Wiredu (décolonisation conceptuelle).
\end{methode}

\begin{exoresolu}[Mobiliser les auteurs : « L'individu et la communauté en Afrique »]
Énoncé : Dans un paragraphe argumenté, utilisez au moins deux auteurs de l'unité pour répondre à la question : la tradition africaine écrase-t-elle l'individu au profit de la communauté ?
\tcblower
\textbf{Solution détaillée.}

\textbf{Introduction de la première référence.} Selon John Mbiti, la formule de la pensée africaine est : « Je suis parce que nous sommes, et parce que nous sommes, je suis. » Explication : cette maxime signifie que l'identité de la personne se construit dans et par la communauté — la famille, le lignage, le village. La communauté n'est donc pas d'abord une contrainte extérieure, mais la condition même de l'épanouissement individuel : au Cameroun, l'enfant porté, éduqué et soutenu par toute la famille élargie en est l'illustration quotidienne.

\textbf{Introduction de la deuxième référence (nuance critique).} Cependant, des philosophes comme Kwasi Wiredu invitent à une « décolonisation conceptuelle » : il faut trier dans l'héritage ce qui mérite d'être conservé. Une communauté qui imposerait tout — mariage forcé, interdiction de toute dissidence — nierait la raison et la liberté de la personne, qui sont des valeurs universelles. Explication : la référence à Wiredu montre que la fidélité à la tradition n'exclut pas l'esprit critique ; c'est même un devoir philosophique.

\textbf{Conclusion du paragraphe.} La tradition africaine n'écrase donc pas nécessairement l'individu : pensée avec Mbiti, elle le fonde ; discutée avec Wiredu, elle doit laisser une place à son autonomie. L'équilibre entre appartenance et liberté est précisément l'un des enjeux majeurs du vivre-ensemble contemporain.
\end{exoresolu}

\begin{attention}[Les 5 erreurs qui coûtent des points au Probatoire / Baccalauréat technique]
\begin{enumerate}
\item \textbf{Décrire la situation sans la philosopher.} Dans les exercices d'application, raconter l'exemple (le conflit, le chef, le repas) sans dégager le problème ni définir la notion (palaver, ubuntu) : c'est du récit, pas de la philosophie. Toujours expliciter le lien notion $\leftrightarrow$ situation.
\item \textbf{Confondre les auteurs et les courants.} Attribuer l'ethnophilosophie à Hountondji ou la négritude à Tempels est une erreur éliminatoire : Tempels et Kagame = ethnophilosophie ; Senghor et Césaire = négritude ; Hountondji, Towa, Wiredu = critique de l'ethnophilosophie ; Mbiti = communautarisme.
\item \textbf{Affirmer sans justifier.} Écrire « la philosophie africaine existe » ou « la gouvernance est mauvaise » sans argument, exemple ni référence. Chaque thèse exige au moins une raison développée.
\item \textbf{Opposer tradition et modernité de façon caricaturale.} Présenter la gouvernance traditionnelle comme parfaite ou comme archaïque est un hors-sujet : le correcteur attend une analyse nuancée (forces et limites des deux modèles, puis leur articulation, par exemple via la décentralisation et les chefferies reconnues).
\item \textbf{Bâcler la conclusion.} Conclusion absente, contradictoire avec le développement, ou relativiste (« chacun a son avis »). La conclusion doit répondre explicitement à la question du sujet et découler des étapes du devoir ; une ouverture maladroite vaut mieux qu'une contradiction.
\end{enumerate}
\end{attention}

\newpage

\coursec{Exercices}

\courssub{Je vérifie mes connaissances}

\begin{exercice}[QCM : repères fondamentaux]
Pour chaque question, une seule réponse est correcte. Recopie la lettre de la bonne réponse.
\begin{enumerate}
\item L'auteur de \emph{L'invention de l'Afrique} est :
\begin{enumerate}
\item[a)] Paulin Hountondji \quad b)] Valentin-Yves Mudimbe \quad c)] Kwasi Wiredu \quad d)] Léopold Sédar Senghor
\end{enumerate}
\item L'ethnophilosophie désigne selon Hountondji :
\begin{enumerate}
\item[a)] la philosophie écrite des universités africaines
\item[b)] un ensemble de croyances collectives présentées à tort comme la philosophie d'un peuple
\item[c)] la philosophie grecque ancienne
\item[d)] la philosophie des droits de l'homme
\end{enumerate}
\item Le principe de l'\cle{Ubuntu} signifie :
\begin{enumerate}
\item[a)] « Je pense, donc je suis »
\item[b)] « Je suis parce que nous sommes »
\item[c)] « L'homme est un loup pour l'homme »
\item[d)] « La fin justifie les moyens »
\end{enumerate}
\item Selon Kwasi Wiredu, la gouvernance traditionnelle akan reposait sur :
\begin{enumerate}
\item[a)] le vote majoritaire \quad b)] le tirage au sort \quad c)] le consensus \quad d)] l'hérédité absolue
\end{enumerate}
\end{enumerate}
\end{exercice}

\begin{exercice}[Vrai ou faux justifié]
Indique si chaque affirmation est vraie ou fausse, puis justifie ta réponse en deux ou trois lignes.
\begin{enumerate}
\item La Négritude et la philosophie africaine sont deux expressions rigoureusement synonymes.
\item La philosophie africaine n'existe que sous forme orale et collective.
\item Le néopatrimonialisme désigne l'usage des ressources publiques de l'État comme s'il s'agissait de biens personnels.
\item La Charte africaine des droits de l'homme et des peuples reconnaît des devoirs à l'individu.
\end{enumerate}
\end{exercice}

\begin{exercice}[Définitions]
Définis en trois à cinq lignes chacune des notions suivantes, en donnant pour chacune un exemple camerounais :
\begin{enumerate}
\item L'\cle{ethnophilosophie}.
\item Le \cle{vivre-ensemble}.
\item La \cle{gouvernance}.
\item Le \cle{consensus}.
\end{enumerate}
\end{exercice}

\begin{exercice}[Questions courtes]
Réponds brièvement à chaque question.
\begin{enumerate}
\item Quelle différence fondamentale distingue la Négritude de la philosophie africaine ?
\item Cite deux devoirs de l'individu énoncés par l'article 29 de la Charte africaine des droits de l'homme et des peuples.
\item Qui est l'auteur de \emph{Sur la philosophie africaine} et quelle thèse centrale y défend-il ?
\item Nomme deux courants ou tendances de la philosophie africaine contemporaine et un auteur associé à chacun.
\end{enumerate}
\end{exercice}

\courssub{Je m'entraîne}

\begin{exercice}[Analyse de situation : la palabre villageoise]
Dans un village de la région de l'Ouest, un conflit foncier oppose deux familles. Le chef convoque une palabre : chaque partie s'exprime longuement, les anciens interviennent, et la discussion se poursuit jusqu'à ce qu'une solution acceptée par tous émerge, sans vote.
\begin{enumerate}
\item À quel modèle de gouvernance traditionnelle cette pratique renvoie-t-elle ?
\item En quoi ce modèle illustre-t-il la thèse de Wiredu sur le consensus ?
\item Relève une force et une limite de ce mode de décision pour la résolution des conflits aujourd'hui.
\end{enumerate}
\end{exercice}

\begin{exercice}[Application : les devoirs du citoyen]
À l'occasion d'une campagne de propreté dans ton quartier à Douala, certains habitants refusent de participer en invoquant leurs droits : « l'État doit tout faire ».
\begin{enumerate}
\item Rappelle deux devoirs de l'individu envers la communauté selon la Charte africaine (article 29).
\item Montre, en t'appuyant sur le principe de l'Ubuntu, pourquoi le vivre-ensemble exige une participation active de chacun.
\item Rédige un court argumentaire (10 lignes maximum) destiné à convaincre tes voisins.
\end{enumerate}
\end{exercice}

\begin{exercice}[Texte à problématiser]
« Une philosophie africaine authentique doit être écrite, critique et individuelle, et non la simple restitution des croyances collectives. »
\begin{enumerate}
\item Identifie le courant de la philosophie africaine auquel cette position appartient et nomme l'auteur qui l'a défendue.
\item Explique ce que cette position reproche à l'ethnophilosophie de Tempels ou de Kagame.
\item Formule la problématique philosophique que soulève ce texte sous la forme d'une question.
\end{enumerate}
\end{exercice}

\begin{exercice}[Tableau comparatif]
Construis un tableau à deux colonnes comparant le \cle{vote majoritaire} et le \cle{consensus} selon les critères suivants : fondement de la légitimité, place de la minorité, rapidité de la décision, risque pour le vivre-ensemble. Conclus en cinq lignes : lequel des deux modes te semble le mieux adapté à une société pluriethnique comme le Cameroun, et pourquoi ?
\end{exercice}

\courssub{Je me prépare au Bac}

\begin{exercice}[Dissertation]
\textbf{Sujet :} Le consensus est-il un mode de gouvernance plus légitime que le vote majoritaire ?

\medskip
Tu traiteras ce sujet en une dissertation complète (introduction problématisée, développement en trois parties, conclusion ouverte). On attend notamment :
\begin{enumerate}
\item Une définition rigoureuse des notions de \cle{consensus}, de \cle{légitimité} et de \cle{vote majoritaire}.
\item La mobilisation de la thèse de Kwasi Wiredu sur la gouvernance consensuelle traditionnelle (akan) et de la critique de Paulin Hountondji.
\item L'examen d'objections : le consensus est-il réalisable dans un État moderne de plusieurs millions de citoyens ? Le vote majoritaire est-il nécessairement source de division ?
\item Une référence à l'expérience camerounaise (diversité culturelle, débats électoraux, dialogue national).
\end{enumerate}
\end{exercice}

\begin{exercice}[Explication de texte : Paulin Hountondji]
\textbf{Texte :} « Par philosophie africaine, j'entends un ensemble de textes, précisément l'ensemble des textes écrits par des Africains et qualifiés par leurs auteurs eux-mêmes de philosophiques. [...] L'ethnophilosophie est une pseudo-philosophie : elle prend pour objet un système de croyances collectives, anonymes, que chacun est censé partager du seul fait d'appartenir à une communauté. Or la philosophie est un discours critique, individuel, qui assume le risque de la contradiction et du débat. La philosophie africaine reste à faire : elle sera une philosophie de combat, engagée dans les luttes de libération de nos peuples, ou elle ne sera pas. » (Paulin Hountondji, \emph{Sur la philosophie africaine}, adapté)

\begin{enumerate}
\item Identifie et reformule en une phrase la thèse centrale du texte.
\item Explique ce que Hountondji reproche à l'ethnophilosophie en relevant deux arguments du texte.
\item Que signifie l'expression « philosophie de combat » ? À quelles luttes pense l'auteur ?
\item Discute la portée de cette thèse pour la recherche universitaire au Cameroun aujourd'hui : la philosophie doit-elle être engagée dans les problèmes de la cité (gouvernance, vivre-ensemble), ou rester un savoir désintéressé ?
\end{enumerate}
\end{exercice}

\begin{exercice}[Cas pratique : conseiller municipal]
Tu es conseiller municipal dans une commune camerounaise de 45\,000 habitants. Le conseil dispose d'un budget exceptionnel de 500 millions de FCFA et doit choisir entre deux projets :
\begin{itemize}
\item \textbf{Projet A :} construire un hôtel de ville moderne, vitrine architecturale de la commune (coût : 480 millions de FCFA) ;
\item \textbf{Projet B :} réhabiliter 10 points d'eau potable dans les villages environnants, où 60\,\% de la population n'a pas d'accès direct à l'eau potable (coût : 450 millions de FCFA).
\end{itemize}
Le maire, élu avec 52\,\% des voix, penche pour le projet A et souhaite trancher par un simple vote du conseil.

\begin{enumerate}
\item Argumente ton choix entre les deux projets en mobilisant l'éthique de la responsabilité (Hans Jonas, Fabien Eboussi Boulaga) : quel projet répond au devoir envers les générations présentes et futures ?
\item Explique en quoi le principe de \cle{subsidiarité} (traiter les problèmes au niveau le plus proche des citoyens) éclaire ce choix.
\item Le maire veut décider par vote majoritaire. Propose, en t'inspirant de la philosophie africaine (consensus, palabre, Ubuntu), une procédure de décision qui associe les populations concernées et renforce le vivre-ensemble.
\item Rédige ta conclusion sous la forme d'une motion de 15 lignes maximum à présenter au conseil municipal.
\end{enumerate}
\end{exercice}

\newpage

\coursec{Corrigés des exercices}

\coursec{La Philosophie africaine}

\courssub{Je vérifie mes connaissances}

\begin{corrige}[Exercice 1 — QCM : repères fondamentaux]
\begin{enumerate}
\item \textbf{Réponse b) : Valentin-Yves Mudimbe.} Dans \emph{L'invention de l'Afrique} (1988), le philosophe congolais Mudimbe montre que « l'Afrique » est en grande partie une construction discursive du discours colonial et des sciences humaines occidentales (anthropologie, mission chrétienne). Paulin Hountondji est l'auteur de \emph{Sur la philosophie africaine} ; Senghor est le théoricien de la Négritude ; Wiredu est le philosophe ghanéen du consensus.
\item \textbf{Réponse b).} Pour Hountondji, l'ethnophilosophie est un ensemble de croyances collectives, anonymes et unanimistes, présentées à tort comme « la » philosophie d'un peuple. Or la philosophie authentique est un discours écrit, critique et individuel.
\item \textbf{Réponse b).} L'\cle{Ubuntu} (mot bantu) signifie « Je suis parce que nous sommes » : l'identité de la personne se construit dans la relation à la communauté. La formule a) est celle de Descartes, c) celle de Hobbes.
\item \textbf{Réponse c) : le consensus.} Selon Kwasi Wiredu, la gouvernance traditionnelle akan (Ghana) reposait sur la délibération prolongée jusqu'à l'accord de tous, et non sur le vote majoritaire qui crée des vainqueurs et des vaincus.
\end{enumerate}
\textbf{Points de vigilance :} ne pas confondre Mudimbe (critique du discours sur l'Afrique) et Hountondji (critique de l'ethnophilosophie) ; ne pas attribuer l'Ubuntu à un auteur unique : c'est un principe de la pensée bantoue, popularisé notamment par Desmond Tutu.
\end{corrige}

\begin{corrige}[Exercice 2 — Vrai ou faux justifié]
\begin{enumerate}
\item \textbf{Faux.} La Négritude (Senghor, Césaire) est un mouvement littéraire et culturel de valorisation de l'identité noire, né dans les années 1930 ; la philosophie africaine est un champ de réflexion critique plus large et postérieur. La Négritude en est une étape historique, non un synonyme.
\item \textbf{Faux.} C'est précisément la thèse combattue par Hountondji : il existe une philosophie africaine écrite, individuelle et critique (Hountondji, Wiredu, Mudimbe, Eboussi Boulaga). Réduire la philosophie africaine à l'oralité collective, c'est retomber dans l'ethnophilosophie.
\item \textbf{Vrai.} Le néopatrimonialisme désigne la confusion entre patrimoine public et patrimoine privé du dirigeant : les ressources de l'État sont gérées comme des biens personnels et distribuées à des clients. C'est une pathologie majeure de la gouvernance, dénoncée par les philosophes africains contemporains.
\item \textbf{Vrai.} La Charte africaine des droits de l'homme et des peuples (1981), dite Charte de Banjul, consacre dans son article 29 des devoirs de l'individu envers sa famille, sa communauté et la société : c'est sa spécificité par rapport aux déclarations occidentales centrées sur les seuls droits.
\end{enumerate}
\end{corrige}

\begin{corrige}[Exercice 3 — Définitions]
\begin{enumerate}
\item \textbf{L'ethnophilosophie} : courant qui présente les croyances, coutumes et visions du monde collectives d'un peuple africain comme sa « philosophie » (Tempels, \emph{La philosophie bantoue} ; Kagame). Critiquée par Hountondji comme pseudo-philosophie anonyme et unanimiste. \emph{Exemple camerounais :} décrire les rites et proverbes beti ou bamiléké comme « la philosophie » de ces peuples relève de l'ethnophilosophie.
\item \textbf{Le vivre-ensemble} : ensemble des principes, valeurs et pratiques (tolérance, solidarité, dialogue, justice) qui permettent à des individus et des groupes différents de cohabiter pacifiquement dans une même société. \emph{Exemple camerounais :} la cohabitation de plus de 250 groupes ethniques et de plusieurs religions au Cameroun, célébrée comme « l'Afrique en miniature ».
\item \textbf{La gouvernance} : manière dont le pouvoir est exercé dans la conduite des affaires publiques ; elle englobe la participation des citoyens, la transparence, la redevabilité (rendre compte) et la bonne gestion des ressources. \emph{Exemple camerounais :} la décentralisation et les conseils municipaux comme cadres de gestion locale des affaires publiques.
\item \textbf{Le consensus} : mode de décision par délibération prolongée jusqu'à ce qu'une solution acceptable par tous émerge, sans vote ni vainqueur désigné. Théorisé par Wiredu à partir de la tradition akan. \emph{Exemple camerounais :} la palabre villageoise, où le chef et les anciens discutent jusqu'à l'accord de toutes les parties.
\end{enumerate}
\end{corrige}

\begin{corrige}[Exercice 4 — Questions courtes]
\begin{enumerate}
\item La Négritude est un mouvement littéraire et identitaire de valorisation de la culture noire (Senghor : « l'émotion est nègre ») ; la philosophie africaine est une démarche rationnelle, critique et écrite qui interroge notamment les conditions de la connaissance et de l'émancipation. La première célèbre une identité, la seconde exerce un examen critique.
\item Selon l'article 29 de la Charte africaine, l'individu a notamment le devoir de \textbf{servir la communauté nationale} en mettant ses capacités physiques et intellectuelles à son service, et le devoir de \textbf{préserver et renforcer la solidarité nationale}, particulièrement lorsque celle-ci est menacée (on peut aussi citer : respecter ses parents, contribuer à la défense de la patrie, préserver les valeurs positives de la civilisation africaine).
\item \emph{Sur la philosophie africaine} (1976) est l'œuvre du Béninois \textbf{Paulin Hountondji}. Il y défend la thèse que la philosophie africaine est l'ensemble des textes écrits par des Africains se revendiquant philosophes, et que l'ethnophilosophie (croyances collectives anonymes) est une pseudo-philosophie : la philosophie est un discours critique, individuel et écrit.
\item Deux tendances possibles parmi : l'\textbf{ethnophilosophie} (Placide Tempels, Alexis Kagame) ; la \textbf{philosophie critique ou « universaliste »} (Paulin Hountondji, Kwasi Wiredu) ; la \textbf{philosophie de la décolonisation du discours} (Valentin-Yves Mudimbe) ; la \textbf{Négritude} (Léopold Sédar Senghor).
\end{enumerate}
\end{corrige}

\courssub{Je m'entraîne}

\begin{corrige}[Exercice 5 — Analyse de situation : la palabre villageoise]
\begin{enumerate}
\item Cette pratique renvoie au modèle de la \cle{gouvernance consensuelle} traditionnelle, fondée sur la palabre : la parole est donnée à tous, la décision naît de la délibération et non du vote.
\item Elle illustre la thèse de \textbf{Wiredu} : dans la tradition akan, on délibère jusqu'à ce que tous les points de vue soient entendus et qu'un accord émerge. Le consensus n'est pas l'unanimité spontanée, mais le résultat d'un travail de persuasion réciproque où chacun accepte de céder sur une partie de ses exigences. Ici, l'absence de vote et la poursuite de la discussion « jusqu'à ce qu'une solution acceptée par tous émerge » correspondent exactement à ce modèle.
\item \textbf{Force :} la solution acceptée par tous est durable, car aucune partie ne se sent vaincue ; elle restaure la paix sociale et le vivre-ensemble (réconciliation des deux familles). \textbf{Limite :} la procédure est longue et difficilement transposable à grande échelle (un État de millions de citoyens ne peut pas palabrer indéfiniment) ; elle peut aussi masquer des rapports de domination (la parole des jeunes, des femmes ou des étrangers au village peut être moins écoutée que celle des anciens).
\end{enumerate}
\end{corrige}

\begin{corrige}[Exercice 6 — Application : les devoirs du citoyen]
\begin{enumerate}
\item Selon l'article 29 de la Charte africaine : le devoir de \textbf{servir la communauté nationale} en mettant ses capacités à son service, et le devoir de \textbf{préserver et renforcer la solidarité nationale}. Participer à la propreté de son quartier relève directement de ces devoirs.
\item Le principe de l'\cle{Ubuntu} (« Je suis parce que nous sommes ») affirme que la personne ne se réalise que dans et par la communauté. Si mon existence même dépend du « nous », alors le bien commun (un quartier propre et sain) est aussi mon bien. Refuser de participer en invoquant ses seuls droits, c'est oublier que droits et devoirs sont solidaires : la Charte africaine les articule l'un à l'autre. Le vivre-ensemble n'est pas un service que l'État livre à des consommateurs, mais une œuvre commune.
\item \textbf{Argumentaire (modèle) :} « Chers voisins, attendre que l'État fasse tout, c'est oublier que nous sommes les premiers bénéficiaires et les premiers responsables de notre cadre de vie. La Charte africaine, qui nous reconnaît des droits, nous assigne aussi des devoirs : servir notre communauté et renforcer la solidarité. Notre tradition le dit depuis toujours : je suis parce que nous sommes. Un quartier propre, c'est moins de maladies pour nos enfants, c'est la dignité pour tous. L'État ne peut pas balayer devant chaque porte ; mais ensemble, en une matinée, nous pouvons transformer notre rue. Donnons chacun un peu de notre temps : le bien commun est notre bien à tous. »
\end{enumerate}
\end{corrige}

\begin{corrige}[Exercice 7 — Texte à problématiser]
\begin{enumerate}
\item Cette position appartient au courant de la \cle{philosophie critique} (ou universaliste), défendue par \textbf{Paulin Hountondji} dans \emph{Sur la philosophie africaine} (1976).
\item Cette position reproche à l'ethnophilosophie de Tempels (\emph{La philosophie bantoue}) et de Kagame trois choses : (i) de confondre des croyances collectives et anonymes avec la philosophie, qui est un discours individuel et assumé par un auteur ; (ii) d'être unanimiste, en prétendant que tous les membres d'un peuple pensent la même chose, ce qui supprime le débat et la contradiction, pourtant constitutifs de la philosophie ; (iii) de rester une description (ethnologie) plutôt qu'une pensée critique qui s'expose à la réfutation.
\item \textbf{Problématique possible :} « La philosophie africaine peut-elle exister autrement que comme discours écrit, critique et individuel ? » ou encore : « Les croyances collectives d'un peuple méritent-elles le nom de philosophie ? »
\end{enumerate}
\end{corrige}

\begin{corrige}[Exercice 8 — Tableau comparatif]
\begin{center}
\begin{tabularx}{\linewidth}{|l|X|X|}\hline
\rowcolor{popblueL} \textbf{Critère} & \textbf{Vote majoritaire} & \textbf{Consensus} \\ \hline
Fondement de la légitimité & Le nombre : la volonté de la majorité (50~\% + 1 voix) fait loi & L'accord de tous après délibération ; chacun a été entendu \\ \hline
Place de la minorité & Vaincue ; elle doit se soumettre à la majorité & Intégrée : ses objections modifient la décision finale \\ \hline
Rapidité de la décision & Rapide : un scrutin tranche & Lente : la palabre peut durer longtemps \\ \hline
Risque pour le vivre-ensemble & Frustration durable des minorités, divisions, « tyrannie de la majorité » & Blocage si un seul refuse ; lenteur paralysante ; consensus de façade imposé par les dominants \\ \hline
\end{tabularx}
\end{center}
\textbf{Conclusion (modèle) :} Dans une société pluriethnique comme le Cameroun, le consensus pur est impraticable à l'échelle nationale (28~millions d'habitants), mais le vote majoritaire seul risque de figer des clivages ethniques ou régionaux : une majorité numérique permanente écraserait des minorités permanentes. La voie la plus adaptée est donc mixte : le vote pour trancher, précédé d'une large délibération (palabre modernisée, dialogue national, concertations) qui associe les groupes concernés et corrige les effets de division du scrutin majoritaire.
\end{corrige}

\courssub{Je me prépare au Bac}

\begin{corrige}[Exercice 9 — Dissertation : le consensus est-il plus légitime que le vote majoritaire ?]
\textbf{Corrigé détaillé (plan et éléments de contenu).}

\textbf{Introduction.} Amorce : toute société doit décider collectivement ; mais comment une décision devient-elle légitime ? Définitions : le \cle{consensus} est un accord obtenu par délibération jusqu'à l'acceptation de tous ; le \cle{vote majoritaire} fait prévaloir la volonté du plus grand nombre ; la \cle{légitimité} est le caractère d'un pouvoir ou d'une décision reconnu comme juste et digne d'obéissance. Annonce du problème : le consensus, hérité des traditions africaines de palabre, semble plus respectueux de chacun ; mais est-il possible et souhaitable dans un État moderne ? \textbf{Problématique :} le consensus est-il un mode de gouvernance plus légitime que le vote majoritaire ? Annonce du plan : I. Oui, le consensus fonde une légitimité supérieure ; II. Mais il se heurte à des limites pratiques et théoriques ; III. Une légitimité démocratique doit articuler les deux.

\textbf{I. La supériorité apparente du consensus.}
\begin{itemize}
\item Thèse de \textbf{Wiredu} : la gouvernance akan traditionnelle reposait sur la délibération consensuelle ; le vote majoritaire, importé d'Occident, institue des vainqueurs et des vaincus, alors que le consensus intègre les objections de chacun. Une décision acceptée par tous est plus légitime qu'une décision subie par $49\,\%$.
\item Fondement anthropologique : l'\cle{Ubuntu} (« je suis parce que nous sommes ») fait de la relation communautaire la condition de la personne ; le consensus exprime cette ontologie relationnelle, le vote majoritaire exprime un individualisme concurrentiel.
\item Enjeu camerounais : dans un pays de plus de 250 groupes ethniques, le consensus préserve le vivre-ensemble là où le scrutin majoritaire peut cristalliser des frustrations identitaires.
\end{itemize}

\textbf{II. Les limites du consensus.}
\begin{itemize}
\item Objection de faisabilité : on ne peut délibérer jusqu'à l'unanimité avec des millions de citoyens ; le consensus exige du temps, une échelle réduite, une culture commune. À l'échelle d'un État moderne, le vote est une nécessité technique.
\item Objection critique (esprit de \textbf{Hountondji}) : le consensus unanimiste peut étouffer la contradiction, qui est le moteur de la pensée critique et du progrès politique ; il peut masquer la domination des anciens, des notables, des hommes sur les femmes et les jeunes — un « consensus » imposé n'est pas légitime.
\item Le vote majoritaire n'est pas nécessairement source de division : encadré par l'alternance, l'État de droit et la protection des minorités, il permet un conflit pacifié et réversible (le vaincu d'aujourd'hui peut gagner demain).
\end{itemize}

\textbf{III. Articuler consensus et vote.}
\begin{itemize}
\item La légitimité démocratique moderne peut combiner les deux : délibérer largement avant de voter (débat public, concertation, dialogue national), voter pour trancher, puis associer les minorités à la mise en œuvre.
\item Référence camerounaise : les grands dialogues et concertations nationales montrent la recherche d'un compromis consensuel au-delà du seul scrutin ; la décentralisation rapproche la décision des citoyens (subsidiarité).
\item Conclusion provisoire : le consensus est un idéal régulateur plus qu'une procédure applicable partout ; le vote est un instrument dont la légitimité dépend de la qualité de la délibération qui le précède.
\end{itemize}

\textbf{Conclusion.} Bilan : le consensus est supérieur en principe (il respecte chaque voix) mais inférieur en pratique à grande échelle ; le vote majoritaire est légitime s'il est précédé et suivi de délibération. Ouverture : la légitimité ne tient-elle pas moins à la procédure qu'à la justice des décisions prises ?

\medskip
\textbf{Barème indicatif (sur 20) :} introduction problématisée avec définitions des trois notions : 4 pts ; développement en trois parties avec thèses, arguments et exemples : 10 pts (dont Wiredu 2 pts, Hountondji 2 pts, référence camerounaise 2 pts) ; examen d'objections : 3 pts ; conclusion ouverte : 1 pt ; langue et cohérence : 2 pts.

\textbf{Points de vigilance :} ne pas opposer caricaturalement « tradition africaine » et « modernité occidentale » ; définir la légitimité dès l'introduction ; citer Wiredu et Hountondji avec exactitude ; éviter la récitation : chaque référence doit servir l'argumentation.
\end{corrige}

\begin{corrige}[Exercice 10 — Explication de texte : Paulin Hountondji]
\begin{enumerate}
\item \textbf{Thèse centrale :} la philosophie africaine authentique est un corpus de textes écrits, individuels et critiques produits par des Africains, et non les croyances collectives anonymes décrites par l'ethnophilosophie ; elle reste à construire comme philosophie engagée dans la libération des peuples.
\item \textbf{Deux reproches à l'ethnophilosophie :} (i) elle prend pour objet « un système de croyances collectives, anonymes », alors que la philosophie est un discours assumé par un individu qui en répond devant la critique ; (ii) elle est unanimiste — « chacun est censé partager » ces croyances « du seul fait d'appartenir à une communauté » —, ce qui supprime « le risque de la contradiction et du débat », c'est-à-dire précisément ce qui fait la philosophie. C'est pourquoi Hountondji la qualifie de « pseudo-philosophie ».
\item « \textbf{Philosophie de combat} » signifie une philosophie engagée, qui met la réflexion critique au service des luttes concrètes des peuples africains : luttes de libération politique (contre les dominations postcoloniales), de développement, de justice sociale et d'émancipation culturelle. L'alternative « ou elle ne sera pas » affirme que la philosophie africaine n'a de sens que si elle s'ancre dans les problèmes réels du continent.
\item \textbf{Discussion (éléments) :} Pour l'engagement : au Cameroun, les philosophes peuvent éclairer les problèmes de gouvernance (néopatrimonialisme, corruption), de vivre-ensemble (diversité ethnique, cohabitation religieuse) et d'éducation ; une recherche coupée de la cité risque la stérilité — c'est l'héritage d'Eboussi Boulaga, qui liait philosophie et exigences de la vie sociale. Pour le désintéressement : la philosophie exige aussi la liberté de penser sans obligation de résultat ; la subordonner à un « combat » risque de la transformer en idéologie et de reproduire l'unanimisme dénoncé. Position nuancée : la recherche universitaire camerounaise peut être rigoureuse et libre dans sa méthode tout en choisissant des objets en prise avec les problèmes de la cité — la rigueur critique est elle-même un service rendu à la démocratie.
\end{enumerate}
\textbf{Barème indicatif (sur 20) :} thèse reformulée : 3 pts ; explication des deux arguments : 6 pts ; sens de « philosophie de combat » : 4 pts ; discussion personnelle argumentée avec exemples camerounais : 5 pts ; langue et organisation : 2 pts.

\textbf{Points de vigilance :} distinguer clairement explication (fidélité au texte) et discussion (prise de position argumentée) ; citer le texte à l'appui ; ne pas réduire Hountondji à un « rejet de la tradition » : il critique une certaine manière de la philosopher, non les cultures africaines.
\end{corrige}

\begin{corrige}[Exercice 11 — Cas pratique : conseiller municipal]
\begin{enumerate}
\item \textbf{Choix argumenté (éthique de la responsabilité).} L'éthique de la responsabilité de \textbf{Hans Jonas} commande d'agir de telle sorte que les effets de l'action soient compatibles avec la permanence d'une vie humaine digne, pour les générations présentes et futures. Chez \textbf{Fabien Eboussi Boulaga}, la philosophie africaine doit répondre aux exigences vitales et à la dignité des personnes concrètes. Or $60\,\%$ de la population n'a pas d'accès direct à l'eau potable : le projet B répond à un besoin vital (santé, dignité, survie des enfants) et constitue un investissement durable pour les générations futures, pour 450 millions de FCFA. Le projet A (480 millions) est une « vitrine » : il sert le prestige du pouvoir plus que le bien commun, et frôle même le néopatrimonialisme symbolique. Le choix responsable est donc le \textbf{projet B}.
\item \textbf{Subsidiarité.} Le principe de subsidiarité veut que les problèmes soient traités au niveau le plus proche des citoyens concernés. Ici, ce sont les villages environnants qui souffrent du manque d'eau : la commune, échelon de proximité, est légitime pour régler ce problème local, et les populations villageoises doivent être associées au choix des sites des points d'eau. Le projet A, lui, ne répond à aucun besoin exprimé par les habitants : il est décidé d'en haut, contre la subsidiarité.
\item \textbf{Procédure de décision inspirée de la philosophie africaine.} Plutôt qu'un vote sec du conseil, on proposera : (i) une \textbf{palabre communale} : réunions publiques dans les villages concernés où chaque habitant peut prendre la parole ; (ii) la recherche du \textbf{consensus} à la manière de Wiredu : délibération du conseil prolongée jusqu'à un accord intégrant les objections des minorités ; (iii) l'esprit de l'\textbf{Ubuntu} : rappeler que le conseil décide pour le « nous » communal, non pour des intérêts particuliers ; (iv) seulement en cas d'échec durable du consensus, un vote en dernier recours, précédé de la publication des arguments. Cette procédure associe les populations et renforce le vivre-ensemble : même les partisans du projet A auront été entendus.
\item \textbf{Motion (modèle, 15 lignes maximum) :}

« Le conseil municipal, réuni en session extraordinaire, considérant que $60\,\%$ de la population de notre commune n'a pas d'accès direct à l'eau potable ; considérant que l'éthique de la responsabilité nous oblige envers les générations présentes et futures ; considérant que le principe de subsidiarité commande de traiter d'abord les besoins vitaux des populations les plus proches ; considérant l'esprit de l'Ubuntu, selon lequel chacun de nous est parce que nous sommes tous ; décide d'affecter 450 millions de FCFA à la réhabilitation de dix points d'eau potable dans les villages environnants ; décide de reporter le projet d'hôtel de ville à un exercice budgétaire ultérieur ; décide enfin d'associer les populations villageoises, par des consultations publiques, au choix des sites et au suivi des travaux. Le prestige d'un bâtiment ne saurait primer sur la santé et la dignité de nos populations : l'eau est la vie, et servir le bien commun est le premier devoir de tout élu. »
\end{enumerate}
\textbf{Barème indicatif (sur 20) :} argumentation éthique (Jonas, Eboussi Boulaga) : 5 pts ; subsidiarité correctement définie et appliquée : 4 pts ; procédure inspirée de consensus/palabre/Ubuntu : 5 pts ; motion rédigée, claire et dans le format : 4 pts ; langue : 2 pts.

\textbf{Points de vigilance :} mobiliser explicitement les auteurs et notions du cours (responsabilité, subsidiarité, Ubuntu, consensus) ; justifier le choix par les chiffres de l'énoncé ($60\,\%$ sans eau potable ; coûts comparables : 450 contre 480 millions) ; la motion doit respecter la forme (« considérant que... ; décide... ») et la limite de 15 lignes.
\end{corrige}

\newpage

\coursec{Fiche bilan}

\begin{popfigure}
\begin{tikzpicture}[mindmap, grow cyclic, every node/.style={concept, text=white, font=\small},
root concept/.append style={concept color=popdark, font=\bfseries},
level 1/.append style={level distance=3.4cm, sibling angle=60},
level 2/.append style={level distance=2.1cm, sibling angle=45, font=\scriptsize}]
\node[root concept]{La Philosophie africaine}
  child[concept color=popblue]{ node {Genèse et historicité}
    child { node {Débat : ethnophilosophie} }
    child { node {Tempels, Hountondji} }
  }
  child[concept color=poporange]{ node {Grandes tendances}
    child { node {Ethnophilosophie} }
    child { node {Sagesse / critique} }
  }
  child[concept color=popgreen]{ node {Vivre-ensemble}
    child { node {Ubuntu : humanité partagée} }
    child { node {Solidarité, palabre} }
  }
  child[concept color=poppurple]{ node {Gouvernance}
    child { node {Modèles traditionnels} }
    child { node {État moderne} }
  }
  child[concept color=poppink]{ node {Défis démocratiques}
    child { node {Élections, justice} }
    child { node {Corruption} }
  }
  child[concept color=popgold]{ node {Enjeux citoyens}
    child { node {Paix, cohésion sociale} }
    child { node {Cameroun : diversité} }
  };
\end{tikzpicture}
\legende{Carte mentale de la leçon : la philosophie africaine, de son histoire à ses thématiques actuelles (vivre-ensemble et gouvernance).}
\end{popfigure}

\begin{bilan}
\textbf{Définitions clés}
\begin{itemize}
  \item \cle{Philosophie africaine} : ensemble des réflexions critiques et rationnelles produites par des Africains sur l'homme, la société et le monde, en dialogue avec leurs cultures et leur histoire.
  \item \cle{Ethnophilosophie} : courant qui identifie la philosophie africaine aux visions du monde collectives et traditionnelles (B. Tempels, \emph{Philosophie bantoue}) ; critiquée par P. Hountondji pour son absence d'esprit critique individuel.
  \item \cle{Ubuntu} : principe sud-africain signifiant « je suis parce que nous sommes » : l'humanité se réalise dans la relation aux autres.
  \item \cle{Vivre-ensemble} : capacité d'une société pluraliste à faire cohabiter harmonieusement des groupes différents autour de valeurs communes (paix, justice, solidarité).
  \item \cle{Gouvernance} : manière d'exercer le pouvoir et d'administrer les affaires publiques ; la \cle{bonne gouvernance} exige transparence, justice, participation et lutte contre la corruption.
  \item \cle{Palabre} : institution traditionnelle africaine de débat public où la parole circule pour résoudre les conflits par le consensus.
\end{itemize}

\textbf{Repères à connaître par cœur}
\begin{center}
\begin{tabularx}{\linewidth}{|l|X|}\hline
\rowcolor{popblueL} \textbf{Auteur / notion} & \textbf{Thèse essentielle} \\ \hline
Tempels & Les Bantous possèdent une philosophie (ontologie de la force vitale). \\ \hline
Hountondji & La philosophie est discours critique individuel, non sagesse collective anonyme. \\ \hline
Ubuntu & L'identité humaine est relationnelle : fondement du vivre-ensemble. \\ \hline
Palabre / consensus & Modèle traditionnel de gouvernance participative et délibérative. \\ \hline
Bonne gouvernance & Transparence, redevabilité, participation citoyenne, état de droit. \\ \hline
\end{tabularx}
\end{center}

\textbf{Méthodes express}
\begin{itemize}
  \item \emph{Dissertation} : définir les termes du sujet dès l'introduction, problématiser (ex. : la tradition africaine est-elle compatible avec la démocratie moderne ?), puis thèse / antithèse / synthèse.
  \item \emph{Application concrète} : toujours illustrer par un exemple camerounais (chefferies traditionnelles, comités de quartier, décentralisation, diversité des 250 groupes ethniques).
  \item \emph{Justifier une réponse} : annoncer la position, donner un argument, appuyer sur un auteur ou un exemple, conclure.
\end{itemize}
\end{bilan}

\begin{aretenir}[Checklist avant le Bac]
\begin{itemize}
  \item[$\square$] Je sais définir la philosophie africaine et situer le débat sur son existence.
  \item[$\square$] Je connais la thèse de Tempels et la critique de Hountondji.
  \item[$\square$] Je distingue ethnophilosophie et philosophie critique.
  \item[$\square$] Je sais expliquer l'Ubuntu et son lien avec le vivre-ensemble.
  \item[$\square$] Je définis la gouvernance et la bonne gouvernance.
  \item[$\square$] Je cite un modèle traditionnel de gouvernance (palabre, chefferie, consensus).
  \item[$\square$] J'identifie les défis démocratiques contemporains (corruption, justice, participation).
  \item[$\square$] Je sais mobiliser un exemple camerounais concret dans une copie.
  \item[$\square$] Je maîtrise la structure d'une dissertation : introduction problématisée, développement en trois parties, conclusion.
  \item[$\square$] Je sais articuler philosophie africaine, vivre-ensemble et gouvernance dans une synthèse.
\end{itemize}
\end{aretenir}

\findecours

\end{document}
